\documentclass[11pt,a4paper]{article}
\usepackage[T1]{fontenc}
\usepackage[utf8]{inputenc}
\usepackage{lmodern}
\usepackage[margin=26mm,headheight=15pt]{geometry}
\usepackage{amsmath,amssymb,amsthm,mathtools}
\usepackage{microtype,booktabs,array,longtable,enumitem,needspace}
\usepackage[dvipsnames]{xcolor}
\usepackage{xurl}
\usepackage[colorlinks=true,linkcolor=MidnightBlue,citecolor=MidnightBlue,urlcolor=MidnightBlue]{hyperref}
\usepackage{fancyhdr}
\pagestyle{fancy}
\fancyhf{}
\fancyhead[L]{\small The logarithmic critical endpoint}
\fancyhead[R]{\small Research article}
\fancyfoot[C]{\thepage}
\renewcommand{\headrulewidth}{0.3pt}
\setlength{\parskip}{4pt}
\setlength{\parindent}{1.2em}
\setlist{itemsep=3pt,topsep=5pt}
\setcounter{tocdepth}{2}
\numberwithin{equation}{section}
\newtheorem{theorem}{Theorem}[section]
\newtheorem{proposition}[theorem]{Proposition}
\newtheorem{lemma}[theorem]{Lemma}
\newtheorem{corollary}[theorem]{Corollary}
\theoremstyle{definition}
\newtheorem{definition}[theorem]{Definition}
\newtheorem{example}[theorem]{Example}
\theoremstyle{remark}
\newtheorem{remark}[theorem]{Remark}
% ed. (2026-09-29): unnumbered environment for ProveIt editorial notes; it does not shift any numbering.
\newtheorem*{ednoteinner}{Editorial note (ProveIt, 2026-09-29)}
\newenvironment{ednote}{\begin{ednoteinner}\sloppy\Urlmuskip=0mu plus 1mu\relax}{\end{ednoteinner}}
\DeclareMathOperator{\Li}{Li}
\DeclareMathOperator{\Log}{Log}
\DeclareMathOperator{\Var}{Var}
\DeclareMathOperator{\Pois}{Pois}
\DeclareMathOperator{\supp}{supp}
\newcommand{\N}{\mathbb N}
\newcommand{\R}{\mathbb R}
\newcommand{\C}{\mathbb C}
\newcommand{\E}{\mathbb E}
\newcommand{\Prob}{\mathbb P}
\newcommand{\dd}{\,\mathrm d}
\newcommand{\bc}{\beta_{\mathrm c}}
\newcommand{\rhoc}{\rho_{\mathrm c}}
\newcommand{\ph}{\phi}
\newcommand{\repo}{\texttt{ProveIt}}
\newcommand{\eps}{\varepsilon}
\hypersetup{pdftitle={The Logarithmic Critical Endpoint: Convergent Lambert Charts, All-Order Sector Laws, and Sharp Action-Cutoff Corrections},pdfauthor={Research draft prepared for Vladimir Reshetnikov},pdfsubject={Transseries, asymptotic inversion, countable exponential feedback, logarithmic criticality, action truncation}}

\begin{document}
\begin{center}
{\small\scshape Transseries / analytic inversion / countable action families}\par
\vspace{12pt}
{\LARGE\bfseries The Logarithmic Critical Endpoint}\par
\vspace{7pt}
{\Large Convergent Lambert Charts, All-Order Sector Laws,\par and Sharp Action-Cutoff Corrections}\par
\vspace{13pt}
{\normalsize Research draft prepared for Vladimir Reshetnikov}\par
{\small 29 September 2026}
\end{center}
\vspace{5pt}
\begin{abstract}
For the positive countable-action equation
\[
 U_\beta(q)=\beta\sum_{j\ge1}a_jq^j e^{jU_\beta(q)},
 \qquad a_j=a j^{-3}\quad\text{outside a finite set},
\]
the critical inverse is neither an analytic square-root fold nor the
fractional Hahn chart associated with a nonintegral power tail. We derive its
logarithmic normal form and construct an exact, convergent two-variable
analytic chart over the negative real branch of Lambert's $W$ function.
An exact Lambert normalization also yields an arbitrary-order, uniformly
additive inverse-logarithmic expansion of the sector coefficients, with
explicit Gaussian--logarithmic Fourier kernels. We then obtain a uniform
finite-coupling refinement of the action-cutoff profile, including its
$1/\sqrt{\log n}$ and $(\log\log n)/\log n$ corrections. Two scales separate:
fluctuations have order $\sqrt{n\log n}$, whereas preserving the $n$th sector
requires and, in the critical window, is equivalent to retaining
$M\gg\sqrt n$ actions. A convergent implicit chart gives all inverse-logarithmic
orders of the finite-action critical-point drift and explains the same cutoff
law from moving analytic folds. A final theorem extends the leading scale
separation to logarithmically weighted cubic tails, including the borderline
$j^{-3}/\log j$ case. Conventional proofs, exact algebraic checks, and
reproducible numerical diagnostics are supplied.
\end{abstract}

\noindent\textbf{Scope and research status.}
The fixed upper endpoint of the earlier countable-action article's Research
Question~4 is treated here; the joint limit in which its exponent approaches
$2$ with $n$ is not. The leading Gaussian/extreme-value separation has
classical predecessors, notably Janson's Example~18.29. The contributions
developed here are the explicit charts, correction formulas, and their
proof-level connection for this model, not a claim to have discovered that
classical mechanism. Publication novelty has not been established by an
exhaustive literature review. No Lean formalization or interval-certified
numerical implementation is claimed.

\clearpage
\begingroup
\small
\setlength{\parskip}{0pt}
\tableofcontents
\endgroup
\clearpage

\section{The research target and its precise boundary}
\label{sec:intro}

\subsection{Why this endpoint is not obtained by substitution}
The repository's transseries program develops polynomial--logarithmic
calculus, inversion around exactly invertible cores, and certified residual
transport. Recent packages also study analytic moving folds and countable
exponential feedback \cite{repo,fold}. The preceding research article
\emph{Beyond Finite-Action Folds} treats an eventually exact tail
$a_j=a j^{-1-\alpha}$ with $1<\alpha<2$ \cite{previous}. At criticality its
normal form begins with $\Gamma(-\alpha)x^\alpha$; the critical chart, the
sector fluctuation scale, and the action budget are organized by this
nonintegral exponent.

The upper endpoint is qualitatively different. The pole of
$\Gamma(-\alpha)$ at $\alpha=2$ cancels a divergent analytic quadratic
coefficient. The surviving expression is $x^2\log(1/x)$, not a finite
multiple of $x^2$. Its inverse needs a logarithmic correction, while its
sector coefficients lie in the infinite-variance Gaussian domain. The
largest relevant individual actions are now smaller than the fluctuation
scale. Substituting $\alpha=2$ into the nonintegral formulas loses all three
features.

Research Question~4 of the earlier article explicitly asks for endpoint
logarithmic charts and cutoff laws. This paper treats the \emph{fixed upper
endpoint}, with finite-prefix perturbations. It also supplies an asymptotic
rate refinement relevant to that article's Question~2, an all-order
finite-fold drift chart relevant to Question~7, and a defined slowly varying
subclass relevant to Question~1. These are scoped resolutions, not a claim
that all of those broader research questions are settled.
% ed. (2026-09-29): an independent second treatment of the same endpoint was filed beside this package (batch 48).
\begin{ednote}
The marginal package \cite{ed:mct}, filed beside this one, treats the same
endpoint (the same model, tail $a j^{-3}$ outside a finite prefix) and was
written independently of this article, also from a library copy of the
earlier article. Its core coincides with this article's: the same
Lambert-$W_{-1}$ core (this article's $(v,T,\eta)$, $h$ and $c$ are its
$(w,t,2r)$, $b$ and $A$), the same chart coefficients and
inverse diagnostics, the same fold functions, and the same critical cutoff
correction, which it states with the sharper error $O(H_n^{-3/2})$. Its
sector constants $\kappa_1,\kappa_2$ are $g_1,g_2$ of
Corollary~\ref{cor:criticalcoeff} rewritten in the normalization
$H_n=\log B_n$ in place of $\ell_n=\log b_n+h$; in particular
$\kappa_1=g_1-h/2=(2-\gamma-\log2-2b)/4$. It is one theorem with two
independent proofs. The Poisson process, the log-weighted tails, the
finite-fold local limit and the exponentiation criterion are only here; its
explicit contraction radii, three-generator logarithmic grid, two-term
minimal budget and finite retention inequality are only there.
\end{ednote}

\subsection{Classical background and the contribution developed here}
The substitution $t=qe^U$ gives $t=q\exp(\beta A(t))$, a classical implicit
function schema. Lagrange inversion and Fourier local-limit methods are
standard tools \cite{flajolet,janson}. In particular, Janson's
Example~18.29 treats cubic power-law weights in the infinite-variance
Gaussian case and explains why conditioning does not change the leading
extreme-value behavior \cite{janson}. The leading scale separation in this
paper must be read in that context, not as an unresolved general
probabilistic phenomenon. The broader distinction between soft and hard
truncations also has an established literature \cite{chakrabarty}.

What is developed explicitly here is a linked analytic and computational
calculus for the action model. The convergent Lambert chart retains the
entire analytic remainder instead of stopping at the leading singular
term. The inverse-logarithmic sector expansion fixes an exact normalization
and gives a formula for every correction. The refined cutoff law keeps
finite coupling, omitted mean, omitted mass, and retained variance
separate. The finite-fold chart identifies a single convergent function
whose Taylor coefficients determine the full inverse-logarithmic drift.
The proofs are included rather than delegated to an unreviewed predecessor.

\begin{center}
\small
\begin{tabular}{>{\raggedright\arraybackslash}p{0.29\textwidth}>{\raggedright\arraybackslash}p{0.60\textwidth}}
\toprule
Result family & Status and contribution in this article\\
\midrule
Implicit substitution and Poisson conditioning & Classical identities, proved here in the exact normalization used later.\\
Leading Gaussian and extreme scales & Classical mechanism; specialized and proved directly for action counts.\\
Lambert--analytic critical chart & Exact convergent normal-form construction, explicit coefficients, and a geometric truncation bound.\\
All-order sector corrections & Explicit Gaussian--logarithmic kernels with uniform additive remainders at every inverse-logarithmic order.\\
Refined cutoff law & Uniform finite-coupling correction with a stated error order, plus an equivalent action-budget criterion.\\
Moving finite folds & Convergent drift and curvature charts with finite-$M$ errors and exponentiation conditions.\\
\bottomrule
\end{tabular}
\end{center}

\subsection{Reading conventions}
All limits have $n\to\infty$, $M\to\infty$, or $\eps\downarrow0$, as
specified. Logarithms of positive numbers are real. Complex logarithms use
the principal determination unless a continuation is explicitly stated.
Uniformity in a parameter on a compact set means that the implicit
constants may depend on that compact set and on the fixed weight family,
but not on $n$, $M$, or the varying parameter. ``All orders'' means one
proved estimate for each fixed truncation order; it does not assert
convergence or optimal truncation of an asymptotic series.

\section{The positive action model and exact coefficient identities}
\label{sec:model}

Fix $a>0$ and a polynomial $P(t)=\sum_{j=1}^Jp_jt^j$ such that
\begin{equation}\label{eq:A}
 A(t)=a\Li_3(t)+P(t)=\sum_{j\ge1}a_jt^j,
 \qquad a_j\ge0,\qquad a_1>0.
\end{equation}
The coefficients $p_j$ themselves need not be nonnegative. Define the finite
prefix moments and the endpoint constants by
\begin{align}
 p^{[k]}&=\sum_{j=1}^J j^kp_j,\qquad
 d_0=a\zeta(3)+p^{[0]},\qquad d_1=a\zeta(2)+p^{[1]},\label{eq:moments}\\
 \bc&=d_1^{-1},\qquad \rho_\beta=e^{-\beta d_0},\qquad
 c_\beta=\beta a,\qquad h=\frac32+\frac{p^{[2]}}a.\label{eq:constants}
\end{align}
Write $c=c_{\bc}$ and $\rhoc=\rho_{\bc}$. Both $d_0$ and $d_1$ are
positive and finite. The genuine second moment $\sum j^2a_j$ diverges;
$p^{[2]}$ is only a finite-prefix quantity.

For $\beta>0$, let
\begin{equation}\label{eq:feedback}
 U_\beta(q)=\beta A(qe^{U_\beta(q)}),\qquad
 u_n(\beta)=[q^n]U_\beta(q).
\end{equation}
The coefficient at degree $n$ is determined by lower coefficients and
$a_1,\ldots,a_n$, so there is a unique formal solution in $q\R[[q]]$.
The analytic implicit function theorem at $(q,U)=(0,0)$ shows that it is
convergent near zero. Its coefficients are positive: the summand with
$a_1>0$ already generates a positive contribution at every degree.

\begin{proposition}[Parametrization and coefficient formula]
\label{prop:exact}
With $t=qe^{U_\beta(q)}$, one has
\begin{equation}\label{eq:param}
 q=t e^{-\beta A(t)},\qquad U_\beta(q)=\beta A(t),\qquad
 u_n(\beta)=\frac1n[t^n]e^{n\beta A(t)}\quad(n\ge1).
\end{equation}
\end{proposition}
\begin{proof}
The first two equalities follow by substitution. Lagrange inversion applied
to $t=q\exp(\beta A(t))$ gives
\[
 [q^n]\beta A(t(q))
 =\frac\beta n[t^{n-1}]A'(t)e^{n\beta A(t)}.
\]
Since $(e^{n\beta A(t)})'=n\beta A'(t)e^{n\beta A(t)}$, coefficient
extraction on the derivative gives the last equality in \eqref{eq:param}.
This also proves the identity purely formally.
\end{proof}

Let $(Z_j)_{j\ge1}$ be independent Poisson variables with means
$n\beta a_j$, and put $S=S_{n,\beta}=\sum_{j\ge1}jZ_j$. The total number
of jumps is finite almost surely because $n\beta\sum a_j<\infty$; also
$\E S=n\beta d_1<\infty$. Its probability generating function is
\[
 \E t^S=\exp\bigl(n\beta(A(t)-d_0)\bigr).
\]
Consequently,
\begin{equation}\label{eq:poisson}
 \boxed{\quad n\rho_\beta^n u_n(\beta)=\Prob(S_{n,\beta}=n).\quad}
\end{equation}

\begin{definition}[Action truncation]
For an integer $M\ge1$, replace $A$ by
$A_M(t)=\sum_{j=1}^M a_jt^j$ in \eqref{eq:feedback}. Denote the resulting
coefficients by $u_n^{(M)}(\beta)$, and set
\[
 Q_M=\sum_{j>M}a_j,\qquad S_M=\sum_{j\le M}jZ_j,\qquad
 R_{n,M}(\beta)=\frac{u_n^{(M)}(\beta)}{u_n(\beta)}.
\]
\end{definition}
Independence and \eqref{eq:poisson} give the exact formula
\begin{equation}\label{eq:cutoff-exact}
 \boxed{\quad
 R_{n,M}=e^{-n\beta Q_M}
 \frac{\Prob(S_M=n)}{\Prob(S=n)}
 =\Prob(Z_j=0\text{ for all }j>M\mid S=n).
 \quad}
\end{equation}
In particular, $0\le R_{n,M}\le1$ and it is nondecreasing in $M$. This
monotonicity will prove the sharp budget criterion without any unproved
uniform extrapolation of a compact-parameter asymptotic formula.

\section{The logarithmic singularity and a convergent Lambert chart}
\label{sec:chart}

\subsection{Deriving the integer-order normal form}
The nonintegral polylogarithm expansion in DLMF~25.12.12 excludes positive
integer orders in its stated form \cite{dlmfpoly}. At order $3$ it is
safer to derive the cancellation directly. In a slit neighborhood of $0$,
\[
 \Li_1(e^{-x})=-\log(1-e^{-x})
 =-\log x+\frac x2-\frac{x^2}{24}+O(x^4).
\]
The difference between $-\log(1-e^{-x})$ and $-\log x$ is analytic for
$|x|<2\pi$: the analytic factor $(1-e^{-x})/x$ has no zero in this disk.
Integrating twice and using the values $\Li_3(1)=\zeta(3)$ and
$\Li_2(1)=\zeta(2)$ gives
\begin{equation}\label{eq:li3}
 \Li_3(e^{-x})=\zeta(3)-\zeta(2)x
 +\frac{x^2}{2}\left(\log\frac1x+\frac32\right)
 +\frac{x^3}{12}-\frac{x^4}{288}+O(x^6).
\end{equation}
The remainder after extracting the displayed logarithm is analytic on
$|x|<2\pi$.

Set
\begin{equation}\label{eq:coordinates}
 x=-\log t,\qquad \eps=\log(\rho_\beta/q),\qquad
 \delta_\beta=1-\beta d_1.
\end{equation}
Equation~\eqref{eq:param} yields the exact local normal form
\begin{equation}\label{eq:normal}
 \boxed{\quad
 \eps=\delta_\beta x+\frac{c_\beta}{2}x^2
     \left(\log\frac1x+h\right)+x^3R_\beta(x),
 \qquad U_\beta=\beta d_0+\eps-x.
 \quad}
\end{equation}
Here $R_\beta$ is analytic for $|x|<2\pi$ and linear in $\beta$, with
\begin{equation}\label{eq:rcoeff}
 R_\beta(0)=\beta\left(\frac a{12}-\frac{p^{[3]}}6\right),\qquad
 R_\beta'(0)=\beta\left(-\frac a{288}+\frac{p^{[4]}}{24}\right).
\end{equation}
At criticality, $\delta_{\bc}=0$. Write $R=R_{\bc}$ and $r_k=[x^k]R(x)$.

\subsection{A general logarithmic inversion theorem}
The following theorem isolates the analytic mechanism from the particular
polylogarithm. Lambert's $W_{-1}$ is the real branch with values at most
$-1$ on $[-e^{-1},0)$ \cite{dlmfw}.

\begin{theorem}[Convergent Lambert--analytic chart]
\label{thm:lambert}
Let $m\ge2$ be an integer, $B>0$, $h\in\R$, and let $R$ be analytic near
zero and real on the real axis. Consider the small positive solution of
\begin{equation}\label{eq:generalnormal}
 \eps=Bx^m\left(\log\frac1x+h\right)+x^{m+1}R(x).
\end{equation}
For all sufficiently small $\eps>0$, define
\begin{equation}\label{eq:generalLambert}
 v=-W_{-1}\left(-\frac{m e^{-mh}\eps}{B}\right),\qquad
 T=\left(\frac{m\eps}{Bv}\right)^{1/m},\qquad \eta=\frac m v.
\end{equation}
There is a function $V(T,\eta)$ holomorphic in a bidisk about $(0,0)$,
with $V(0,0)=1$, such that the solution is exactly
\begin{equation}\label{eq:Veq}
 x=TV(T,\eta),\qquad
 V^m(1-\eta\log V)+\frac\eta B T V^{m+1}R(TV)=1.
\end{equation}
Moreover, $V(T,0)=V(0,\eta)=1$. Thus $V-1$ is divisible by $T\eta$ in
the local ring of convergent power series. The representation also gives
the analytic continuation of this branch in a sufficiently small sector
about the positive $\eps$ axis, with compatible branches in
\eqref{eq:generalLambert}.
\end{theorem}
\begin{proof}
The defining relation for $W$ gives
$v e^{-v}=m e^{-mh}\eps/B$. It follows that
\[
 \log(1/T)+h=v/m=1/\eta,\qquad
 BT^m\bigl(\log(1/T)+h\bigr)=\eps.
\]
Substituting $x=TV$ into \eqref{eq:generalnormal} and dividing by this
last identity gives \eqref{eq:Veq}. Its left side minus $1$ is holomorphic
near $(V,T,\eta)=(1,0,0)$, taking the branch of $\log V$ near $1$.
The derivative with respect to $V$ at that point is $m\ne0$. The analytic
implicit function theorem therefore gives a unique $V$ near $1$.
At $T=0$, $V=1$ is the locally unique solution; the same holds at
$\eta=0$. Divisibility follows from the convergent double Taylor series.

For positive $\eps$ tending to zero, $v\to\infty$, $T\to0$ and
$\eta\to0$, so the constructed real $x$ is positive. The derivative of
the right side of \eqref{eq:generalnormal} is
\[
 Bx^{m-1}\left(m\log\frac1x+mh-1\right)+O(x^m)>0
\]
for sufficiently small $x>0$. Hence it is the unique small positive
solution. The branch functions and the bidisk solution admit compatible
analytic continuation to a small sector about this positive ray.
\end{proof}

\begin{corollary}[Critical chart for cubic action weights]
\label{cor:criticalchart}
In \eqref{eq:normal} at $\beta=\bc$, put
\begin{equation}\label{eq:criticalLambert}
 v=-W_{-1}\left(-\frac{4e^{-2h}\eps}{c}\right),\qquad
 T=2\sqrt{\frac{\eps}{cv}},\qquad \eta=\frac2v.
\end{equation}
Then $x=TV(T,\eta)$ is a convergent chart and
\begin{align}
 V(T,\eta)&=1+v_1(\eta)T+v_2(\eta)T^2+O(T^3),\label{eq:Vexp}\\
 v_1(\eta)&=-\frac{2\eta r_0}{c(2-\eta)},\label{eq:vone}\\
 v_2(\eta)&=-\frac{2\eta r_1}{c(2-\eta)}
 +\frac{2\eta^2(10-3\eta)r_0^2}{c^2(2-\eta)^3}.\label{eq:vtwo}
\end{align}
The $O(T^3)$ is uniform for $\eta$ in a sufficiently small closed disk.
At each fixed order in $T$, its coefficient is a rational function of
$\eta$ whose possible poles are at $\eta=2$.
\end{corollary}
\begin{proof}
Take $m=2$ and $B=c/2$ in Theorem~\ref{thm:lambert}. At $T=0$ the
linearized coefficient in $V-1$ is $2-\eta$. Comparing the coefficients
of $T$ and $T^2$ gives \eqref{eq:vone}--\eqref{eq:vtwo}. At every higher
order the new unknown appears linearly with that same coefficient;
the remaining expression is polynomial in previous coefficients and
finitely many $r_k$, with polynomial dependence on $\eta$. Induction
proves the rational-function assertion.
\end{proof}

\begin{proposition}[Geometric chart truncation]
\label{prop:geometric}
Choose $r,s,C>0$ so that $V$ is holomorphic on a neighborhood of the
closed bidisk $|T|\le r$, $|\eta|\le s$ and $|V|\le C$ there. Write
$V=\sum_{j,k\ge0}v_{jk}T^j\eta^k$. If $q_T=|T|/r<1$ and
$q_\eta=|\eta|/s<1$, then
\begin{equation}\label{eq:geometric}
 \left|x-T\sum_{j=0}^{J}\sum_{k=0}^{K}v_{jk}T^j\eta^k\right|
 \le |T|C\frac{q_T^{J+1}+q_\eta^{K+1}}
 {(1-q_T)(1-q_\eta)}.
\end{equation}
\end{proposition}
\begin{proof}
Cauchy's inequalities give $|v_{jk}|\le Cr^{-j}s^{-k}$. Sum the two
geometric tails for $j>J$ and for $k>K$; counting their intersection
twice only increases the bound.
\end{proof}
The theorem provides the existence of such a bidisk. Obtaining numerical
values of $r,s,C$ from a chosen bound on $R$ is a separate local
majorant calculation; the supplied numerical program does not certify
those values.

\subsection{What the chart says about ramification and radius}
\begin{proposition}[Logarithmic obstruction to finite Puiseux ramification]
\label{prop:ramification}
At criticality,
\begin{equation}\label{eq:leadingx}
 x\sim2\sqrt{\frac{\eps}{c\log(1/\eps)}}.
\end{equation}
There is no positive integer $N$ for which the small critical inverse
$x(\eps)$ is represented by a nonzero convergent power series in
$\eps^{1/N}$ near $0$. The radius of convergence of $U_{\bc}(q)$ at
$q=0$ is $\rhoc$.
\end{proposition}
\begin{proof}
The defining equation $v-\log v=\log(c e^{2h}/(4\eps))$ gives
$v\sim\log(1/\eps)$. Since $V\to1$, \eqref{eq:leadingx} follows.
A convergent Puiseux series vanishing at zero has a first nonzero term
$C\eps^{k/N}$. The logarithmic order in \eqref{eq:leadingx} forces
$k/N=1/2$, but $x/\sqrt\eps\to0$, contradicting $C\ne0$.

For $0<r<1$, positivity and \eqref{eq:param} give
$u_n(\bc)\le e^{n\bc A(r)}/(nr^n)$. Letting $r\uparrow1$ shows that
the radius is at least $\rhoc$. On $0<t<1$,
\[
 \frac{\dd}{\dd t}\bigl(t e^{-\bc A(t)}\bigr)
 =e^{-\bc A(t)}\bigl(1-\bc tA'(t)\bigr)>0,
\]
because $tA'(t)<d_1$. The inverse branch reaches $t=1$ when
$q=\rhoc$. If $U_{\bc}$ were holomorphic there, then
$x=\bc d_0+\eps-U_{\bc}(q)$ would be holomorphic in $\eps$ at $0$,
contrary to the first part. Thus the radius is exactly $\rhoc$.
\end{proof}

The obstruction is logarithmic, not an analytic critical point with a
new integer multiplicity. The critical inverse is nonetheless described
by an \emph{actually convergent} analytic function of two explicit small
Lambert coordinates. Expanding $W_{-1}$ itself produces the familiar
nested logarithms; that optional further expansion is not needed for
\eqref{eq:geometric}.

\section{Lambert normalization of the sector scale}
\label{sec:scale}

Let $\beta$ range over a fixed compact subinterval of $(0,\infty)$.
For sufficiently large $n$, define the large positive solution of
\begin{equation}\label{eq:bn}
 b_n^2=n c_\beta(\log b_n+h),\qquad
 \ell_n=\log b_n+h,\qquad e_n=\sqrt{n c_\beta}.
\end{equation}
The letter $e_n$ denotes the \emph{extreme-action scale}, not Euler's
constant or an exponential function. These quantities have the exact
form
\begin{equation}\label{eq:ellLambert}
 \boxed{\quad
 \ell_n=-\frac12W_{-1}\left(-\frac{2e^{-2h}}{n c_\beta}\right),
 \qquad b_n=e_n\sqrt{\ell_n}.
 \quad}
\end{equation}
Indeed, $e^{2\ell_n-2h}=n c_\beta\ell_n$ is equivalent to the
Lambert equation. In particular,
\begin{equation}\label{eq:scalesleading}
 \ell_n\sim\frac12\log n,\qquad
 b_n\sim\sqrt{\frac{n c_\beta\log n}{2}},\qquad
 \frac{e_n}{b_n}=\ell_n^{-1/2}\longrightarrow0.
\end{equation}
The constants $h$ and $c_\beta$ are retained in \eqref{eq:ellLambert};
replacing them prematurely by a leading logarithm changes the correction
coefficients below.

For a general integer $k$, define
\begin{equation}\label{eq:lambda}
 \lambda=\frac{k-n\beta d_1}{b_n}.
\end{equation}
For the desired sector coefficient, $k=n$ and
$\lambda_n=n(1-\beta d_1)/b_n$. Thus a compact critical window
$|\lambda_n|\le R$ has coupling width
\begin{equation}\label{eq:couplingwidth}
 \beta-\bc=O\left(\sqrt{\frac{\log n}{n}}\right).
\end{equation}
We shall allow $\beta=\beta_n$ to vary within this window. The
all-order local-limit theorem itself is stronger: it is uniformly
additive for every integer $k$ and for $\beta$ in a fixed compact
positive interval.

\section{An all-order Gaussian--logarithmic sector expansion}
\label{sec:allorders}

\subsection{The correction kernels}
Let $\ph(\lambda)=(2\pi)^{-1/2}e^{-\lambda^2/2}$ and define, for
$j\ge0$,
\begin{equation}\label{eq:Gj}
 G_j(\lambda)=\frac{1}{2\pi\,2^j j!}
 \int_{\R}e^{-s^2/2-i\lambda s}s^{2j}
      [\Log(-is)]^j\dd s.
\end{equation}
At $s=0$ the integrand is interpreted by its integrable limit when
$j\ge1$. The branch is
$\Log(-is)=\log|s|-i(\pi/2)\operatorname{sgn}(s)$.
The integrals define entire functions of $\lambda$: on each compact
complex set, the Gaussian dominates every derivative of the integrand.
They are real on the real axis by conjugation symmetry, and $G_0=\ph$.

\begin{theorem}[Uniform all-order local limit]
\label{thm:allorders}
For each integer $K\ge0$, uniformly for $\beta$ in a compact subinterval
of $(0,\infty)$ and all integers $k$,
\begin{equation}\label{eq:localexp}
 \boxed{\quad
 b_n\Prob(S_{n,\beta}=k)
 =\sum_{j=0}^{K}\frac{G_j(\lambda)}{\ell_n^j}
       +O(\ell_n^{-K-1}),
 \qquad \lambda=\frac{k-n\beta d_1}{b_n}.
 \quad}
\end{equation}
Consequently, for $k=n$,
\begin{equation}\label{eq:sectorexp}
 u_n(\beta)=\frac{\rho_\beta^{-n}}{n b_n}
 \left\{\sum_{j=0}^{K}\frac{G_j(\lambda_n)}{\ell_n^j}
       +O(\ell_n^{-K-1})\right\}.
\end{equation}
These estimates are additive. Dividing by $\ph(\lambda)$ gives a uniform
relative expansion only on compact real $\lambda$ intervals.
\end{theorem}

\subsection{Fourier damping and a complete proof}
\begin{lemma}[Aperiodic logarithmic damping]
\label{lem:damping}
There are positive constants $C_1,C_2,\theta_0$ such that
\begin{align}
 \sum_{j\ge1}a_j(1-\cos j\theta)
 &\ge C_1\theta^2\log(1/|\theta|)
       &&(0<|\theta|\le\theta_0),\label{eq:damping-small}\\
 \sum_{j\ge1}a_j(1-\cos j\theta)
 &\ge C_2 &&(\theta_0\le|\theta|\le\pi).
 \label{eq:damping-away}
\end{align}
For the truncated sum, the first lower bound is replaced, after reducing
$\theta_0$ and taking $M$ large, by
\begin{equation}\label{eq:damping-truncated}
 \sum_{j\le M}a_j(1-\cos j\theta)
 \ge C_1\theta^2\log\bigl(\min(M,1/|\theta|)\bigr).
\end{equation}
The away-from-zero estimate remains valid for $M\ge1$.
\end{lemma}
\begin{proof}
For $|j\theta|\le1$, $1-\cos(j\theta)\ge c_0j^2\theta^2$.
Sum over the unaltered tail indices
$J<j\le\min(M,1/|\theta|)$ and use
$\sum_{J<j\le L}j^{-1}\ge\frac12\log L$ for large $L$.
This proves the small-angle bounds with fixed smaller constants.
For $\theta_0\le|\theta|\le\pi$, the term
$a_1(1-\cos\theta)$ has a strictly positive minimum.
\end{proof}

\begin{proof}[Proof of Theorem~\ref{thm:allorders}]
Fourier inversion gives
\begin{equation}\label{eq:Fourier}
 b_n\Prob(S=k)=\frac1{2\pi}
 \int_{-\pi b_n}^{\pi b_n}
 e^{-i\lambda s}\exp\{\Psi_n(s)\}\dd s,
\end{equation}
where
\[
 \Psi_n(s)=n\beta\left(A(e^{is/b_n})-d_0-i d_1s/b_n\right).
\]
From \eqref{eq:normal}, with $x=-is/b_n$, and the defining equality
$b_n^2=n c_\beta\ell_n$, we obtain
\begin{equation}\label{eq:char-normal}
 \Psi_n(s)=-\frac{s^2}{2}
   +\frac{s^2}{2\ell_n}\Log(-is)
   +O\left(\frac{n|s|^3}{b_n^3}\right)
\end{equation}
uniformly when $|s|/b_n$ is sufficiently small. The analytic remainder
in \eqref{eq:normal} is uniformly bounded as $\beta$ varies in the
specified compact interval.

We first justify localization. For $|s|\le b_n^{1/2}$, the logarithm in
\eqref{eq:damping-small} is at least $(\log b_n)/2$, so the modulus of
$e^{\Psi_n(s)}$ is bounded by $e^{-C s^2}$. For
$b_n^{-1/2}\le|\theta|\le\theta_0$, the function
$\theta^2\log(1/|\theta|)$ is increasing after choosing
$\theta_0$ small. The damping exponent is at least a positive multiple
of $b_n$. For $\theta_0\le|\theta|\le\pi$ it is at least a positive
multiple of $n$. Their contributions to \eqref{eq:Fourier}, including
the factor $b_n$, are smaller than every power of $1/\ell_n$.

Fix $K$ and choose $D$ large. On
$|s|\le S_n=\sqrt{D\log\ell_n}$, the last term of
\eqref{eq:char-normal} is negligible to every fixed inverse-logarithmic
order, since $n/b_n^3=O(n^{-1/2}(\log n)^{-3/2})$.
Expand
\[
 \exp\left(\frac{s^2}{2\ell_n}\Log(-is)\right)
 =\sum_{j=0}^K\frac{s^{2j}[\Log(-is)]^j}
 {2^j j!\ell_n^j}+\mathcal E_{K,n}(s).
\]
The integral of $e^{-s^2/2}|\mathcal E_{K,n}(s)|$ on this interval is
$O(\ell_n^{-K-1})$. To see this without a growing logarithmic loss,
use the exponential Taylor remainder. Its extra factor
\[
 \exp\bigl(O(s^2|\Log(-is)|/\ell_n)\bigr)
\]
can be absorbed into $e^{-s^2/4}$; the resulting integral of $|s|^{2K+2}(1+|\log|s||)^{K+1}e^{-s^2/4}$ is finite
and independent of $n$. The analytic remainder in
\eqref{eq:char-normal} is handled in the same way.

The part $S_n<|s|\le b_n^{1/2}$ is $O(\ell_n^{-K-1})$ by the
Gaussian damping, after increasing $D$. The tails of the individual
kernel integrals satisfy the same estimate. This gives
\eqref{eq:localexp}. Every bound is independent of $\lambda$, because
$|e^{-i\lambda s}|=1$ for real $\lambda$. Finally use
\eqref{eq:poisson} to obtain \eqref{eq:sectorexp}.
\end{proof}

\subsection{Explicit critical coefficients}
\begin{corollary}[First critical corrections and a formula at every order]
\label{cor:criticalcoeff}
At $\beta=\bc$,
\begin{equation}\label{eq:criticalcoeff}
 u_n(\bc)=\frac{\rhoc^{-n}}{\sqrt{2\pi}\,n b_n}
 \left(1+\frac{g_1}{\ell_n}+\frac{g_2}{\ell_n^2}
       +\cdots+\frac{g_K}{\ell_n^K}+O(\ell_n^{-K-1})\right),
\end{equation}
where, with $\gamma$ denoting Euler's constant,
\begin{align}
 g_1&=\frac{2-\gamma-\log2}{4},\label{eq:g1}\\
 g_2&=\frac3{32}\left[
    \left(\frac83-\gamma-\log2\right)^2
                -\frac{\pi^2}{2}-\frac{40}{9}\right].\label{eq:g2}
\end{align}
For every $j\ge0$, define
\[
 M_j(z)=\frac{2^{j+z/2}\Gamma(j+(z+1)/2)}{\sqrt\pi}.
\]
Then
\begin{equation}\label{eq:gj-moments}
 g_j=\frac1{2^j j!}\Re\left[
       \left(\frac{\dd}{\dd z}-\frac{i\pi}{2}\right)^j M_j(z)
                            \right]_{z=0}.
\end{equation}
In particular,
\begin{equation}\label{eq:leadingsector}
 u_n(\bc)\sim\frac{\rhoc^{-n}}
 {\sqrt{\pi c}\,n^{3/2}\sqrt{\log n}}.
\end{equation}
\end{corollary}
\begin{proof}
At criticality $\lambda_n=0$. Divide \eqref{eq:Gj} by $\ph(0)$ and
use the even symmetry to write it as
$(2^j j!)^{-1}\E[|Z|^{2j}\Re(\log|Z|-i\pi/2)^j]$, where $Z$ is
standard normal. Its Mellin moment is $M_j(z)$; differentiation under
the integral is justified by Gaussian decay and integrability at zero.
This proves \eqref{eq:gj-moments}. The elementary half-integer gamma
recurrences yield \eqref{eq:g1} and \eqref{eq:g2}. Finally use
\eqref{eq:scalesleading}.
\end{proof}

The finite prefix changes $d_0,d_1,h$, and therefore changes the
exponential factor and exact Lambert scale. After that normalization,
the kernels $G_j$ do not depend on the prefix. Its remaining analytic
information, carried by $R_\beta$, contributes terms of order
$n/b_n^3$ and smaller in the Fourier exponent. These are beyond all
fixed inverse-logarithmic orders, but they need not be absent. This is
a precise boundary of the universality statement.

\section{Sharp action budgets and finite-coupling cutoff corrections}
\label{sec:cutoff}

Throughout this section, $\beta=\beta_n$ lies in a fixed compact positive
interval, and the sector coordinate
$\lambda_n=n(1-\beta_n d_1)/b_n$ lies in a fixed compact real interval.
In particular, $\beta_n\to\bc$. Use the scales in
\eqref{eq:bn} with this varying $\beta_n$, and put
\begin{equation}\label{eq:sn}
 s_n=\frac{M}{e_n}.
\end{equation}
The first theorem gives the leading law. Its refinement then quantifies
why convergence to that law is slow.

\begin{theorem}[Cutoff profile and necessary-and-sufficient budget]
\label{thm:budget}
Uniformly when $s_n$ and $\lambda_n$ range over compact subsets of
$(0,\infty)$ and $\R$, respectively,
\begin{equation}\label{eq:leadingcutoff}
 R_{n,M}(\beta_n)=\exp\left(-\frac1{2s_n^2}\right)+o(1).
\end{equation}
For arbitrary integer cutoffs $M=M_n\ge1$, subject to the same compact
critical-window condition,
\begin{equation}\label{eq:budgetiff}
 \boxed{\quad R_{n,M_n}(\beta_n)\longrightarrow1
       \quad\Longleftrightarrow\quad M_n/e_n\longrightarrow\infty.\quad}
\end{equation}
If $M_n/e_n\to0$, the retained fraction tends to zero. Thus the sharp
budget scale is $e_n\asymp\sqrt n$, not $b_n\asymp\sqrt{n\log n}$.
\end{theorem}

\subsection{A quantitative local limit for the retained actions}
The exact omitted mass, omitted mean on the full fluctuation scale, and
retained variance on that scale are
\begin{equation}\label{eq:nukappav}
 \nu_{n,M}=n\beta Q_M,\qquad
 \kappa_{n,M}=\frac{n\beta}{b_n}\sum_{j>M}j a_j,\qquad
 v_{n,M}=\frac{n\beta}{b_n^2}\sum_{j\le M}j^2a_j.
\end{equation}
They must not be confused with the exponentially tilted variance of a
finite fold in Section~\ref{sec:folds}. Let
$\ph_v(z)=(2\pi v)^{-1/2}\exp(-z^2/(2v))$.

\begin{lemma}[Truncated Gaussian approximation with a rate]
\label{lem:truncated}
If $s_n=M/e_n$ ranges over a compact subset of $(0,\infty)$, then,
uniformly for all integers $k$,
\begin{equation}\label{eq:truncatedLLT}
 b_n\Prob(S_M=k)
 =\ph_{v_{n,M}}\left(\frac{k-n\beta d_1}{b_n}
                         +\kappa_{n,M}\right)
       +O(\ell_n^{-3/2}).
\end{equation}
\end{lemma}
\begin{proof}
Write $\mu_M=n\beta\sum_{j\le M}j a_j$. The centered characteristic
exponent, at frequency $s/b_n$, equals
\[
 n\beta\sum_{j\le M}a_j
     (e^{ijs/b_n}-1-ijs/b_n)
   =-\frac{v_{n,M}s^2}{2}+\mathcal R_n(s),
\]
with
\begin{equation}\label{eq:thirdcumulant}
 |\mathcal R_n(s)|\le
 \frac{n\beta |s|^3}{6b_n^3}\sum_{j\le M}j^3a_j
 =O(\ell_n^{-3/2}|s|^3).
\end{equation}
Here $\sum_{j\le M}j^3a_j=aM+O(1)$ and $M\asymp e_n$.
Also
\[
 v_{n,M}=\frac{H_M+p^{[2]}/a}{\ell_n}\longrightarrow1,
 \qquad H_M=\sum_{j=1}^M\frac1j,
\]
because $\log M=\log b_n-\tfrac12\log\ell_n+O(1)$.

For $|s|\le\sqrt{D\log\ell_n}$, integrate the difference between the
true centered characteristic function and $e^{-v_{n,M}s^2/2}$.
Using \eqref{eq:thirdcumulant}, the Gaussian damping, and the identity
$e^z-e^w=(z-w)\int_0^1 e^{(1-t)w+tz}\dd t$, its integral is
$O(\ell_n^{-3/2})$. The remainder of the Fourier interval is negligible
to this order. Indeed, for $|\theta|\le M^{-1/2}$,
\eqref{eq:damping-truncated} gives a Gaussian bound on the $b_n$ scale,
since $\log\min(M,1/|\theta|)\ge\tfrac12\log M\asymp\ell_n$.
For $M^{-1/2}\le|\theta|\le\theta_0$ the exponent is at least a
constant times $n\log M/M$, which is much larger than $\log b_n$.
The remaining interval is controlled by $a_1>0$.
Fourier inversion now gives \eqref{eq:truncatedLLT}; the shift is
$(k-\mu_M)/b_n=(k-n\beta d_1)/b_n+\kappa_{n,M}$.
The absolute integral bound is independent of $k$.
\end{proof}

\begin{theorem}[Uniform refined cutoff formula]
\label{thm:refined}
Put $g(\lambda)=G_1(\lambda)/\ph(\lambda)$. Uniformly for $s_n$ in a
compact positive interval and $\lambda_n$ in a compact real interval,
\begin{equation}\label{eq:refinedexact}
 \boxed{\begin{aligned}
 \log R_{n,M}={}&-\nu_{n,M}-\frac12\log v_{n,M}
       +\frac{\lambda_n^2}{2}
       -\frac{(\lambda_n+\kappa_{n,M})^2}{2v_{n,M}}\\
 &\hspace{10mm}-\frac{g(\lambda_n)}{\ell_n}
       +O(\ell_n^{-3/2}).
 \end{aligned}}
\end{equation}
For $M>J$, define
\begin{equation}\label{eq:D}
 D_n=\frac12\log\ell_n-\log s_n-\gamma+\frac32.
\end{equation}
Then the equivalent expanded form is
\begin{equation}\label{eq:refinedexpanded}
 \boxed{\begin{aligned}
 \log R_{n,M}={}&-\frac1{2s_n^2}
       -\frac{\lambda_n}{s_n\sqrt{\ell_n}}\\
 &+\frac1{\ell_n}
   \left(\frac{1-\lambda_n^2}{2}D_n
        -\frac1{2s_n^2}-g(\lambda_n)\right)
   +O\left(\frac{\log\ell_n}{\ell_n^{3/2}}\right).
 \end{aligned}}
\end{equation}
\end{theorem}
\begin{proof}
Theorem~\ref{thm:allorders} with $K=1$ gives
\[
 b_n\Prob(S=n)=\ph(\lambda_n)
       \left(1+\frac{g(\lambda_n)}{\ell_n}
                   +O(\ell_n^{-2})\right).
\]
Lemma~\ref{lem:truncated} gives the corresponding numerator. On the
specified compact sets the Gaussian values stay bounded away from zero,
because $v_{n,M}\to1$ and $\kappa_{n,M}\to0$. Taking logarithms in
\eqref{eq:cutoff-exact} proves \eqref{eq:refinedexact}.

For $M>J$, integral comparison or Euler summation gives
\begin{align*}
 \sum_{j>M}a_j&=a\left(\frac1{2M^2}+O(M^{-3})\right),\\
 \sum_{j>M}j a_j&=a\left(\frac1M+O(M^{-2})\right),\\
 \sum_{j\le M}j^2a_j&=a(\log M+\gamma)+p^{[2]}+O(M^{-1}).
\end{align*}
Using $M=s_ne_n$, $b_n=e_n\sqrt{\ell_n}$, and
$\ell_n=\log b_n+3/2+p^{[2]}/a$ yields
\begin{align}
 \nu_{n,M}&=\frac1{2s_n^2}+O(n^{-1/2}),\label{eq:tail-exp}\\
 \kappa_{n,M}&=\frac1{s_n\sqrt{\ell_n}}
                    +O(n^{-1/2}\ell_n^{-1/2}),\label{eq:mean-exp}\\
 v_{n,M}&=1-\frac{D_n}{\ell_n}
                    +O(n^{-1/2}\ell_n^{-1}).\label{eq:var-exp}
\end{align}
In particular the explicit $p^{[2]}$ cancels in $D_n$. Expand the
logarithm and reciprocal of $v_{n,M}$ in \eqref{eq:refinedexact}.
The first cross term from $\kappa_{n,M}$ is
$-\lambda_n/(s_n\sqrt{\ell_n})$; the square of that shift contributes
$-1/(2s_n^2\ell_n)$. The linear variance correction is
$(1-\lambda_n^2)D_n/(2\ell_n)$. The unused products are bounded by
$O(D_n\ell_n^{-3/2}+D_n^2\ell_n^{-2})$, which is
$O((\log\ell_n)\ell_n^{-3/2})$ for large $n$.
This proves \eqref{eq:refinedexpanded}.
\end{proof}

At exact criticality, $\lambda_n=0$, so \eqref{eq:refinedexpanded}
simplifies to
\begin{equation}\label{eq:criticalcutoffcorrection}
 \log R_{n,M}=-\frac1{2s_n^2}
 +\frac1{\ell_n}\left[
       \frac14\log\ell_n-\frac12\log s_n
       +\frac{1-\gamma+\log2}{4}-\frac1{2s_n^2}\right]
 +O\left(\frac{\log\ell_n}{\ell_n^{3/2}}\right).
\end{equation}
The $(\log\ell_n)/\ell_n$ term explains a particularly slow approach
to the leading cutoff profile. Away from $\lambda_n=0$, the larger
$1/\sqrt{\ell_n}$ correction records the omitted mean. It cannot be
recovered from an omitted-intensity estimate alone.

\begin{proof}[Proof of Theorem~\ref{thm:budget}]
The leading formula follows immediately from
Theorem~\ref{thm:refined}. To extend it to arbitrary cutoffs, use
monotonicity in \eqref{eq:cutoff-exact}. If $M_n/e_n\to\infty$, then
for every fixed $L>0$,
\[
 \liminf_{n\to\infty}R_{n,M_n}
 \ge\lim_{n\to\infty}R_{n,\lfloor Le_n\rfloor}
 =e^{-1/(2L^2)}.
\]
Let $L\to\infty$ and use $R_{n,M_n}\le1$.
Conversely, if $M_n/e_n$ does not tend to infinity, some subsequence
has $M_n/e_n\le L$ for a finite $L$. Along it, monotonicity bounds the
limsup by $e^{-1/(2L^2)}<1$ (enlarging $L$ slightly if necessary).
Finally, if $M_n/e_n\to0$, bound above by the profile at
$\lfloor Le_n\rfloor$ for every fixed $L>0$ and then let $L\downarrow0$.
\end{proof}

\subsection{Action selection and the limit of the guarantee}
For a desired asymptotic retained fraction $r\in(0,1)$, a natural budget
is
\begin{equation}\label{eq:selection}
 M_n\sim\frac{\sqrt{n\beta_n a}}{\sqrt{2\log(1/r)}}.
\end{equation}
For example, as the allowed loss $1-r$ tends to zero after the asymptotic
limit, this factor behaves like $1/\sqrt{2(1-r)}$. Formula
\eqref{eq:refinedexact} gives a more informative finite-$n$ prediction,
because its omitted sums can be evaluated exactly or tightly bounded.
It is not, as presently implemented, a certified finite-$n$ relative-error
algorithm: explicit constants and threshold values in the $O$ term would
be required for that stronger deliverable.

\section{The conditional point process of large actions}
\label{sec:pointprocess}

The cutoff profile has a transparent interpretation. Although a large
individual action has size of order $e_n$, its contribution divided by
$b_n$ tends to zero. Conditioning on a fluctuation in the $b_n$ window
therefore leaves the leading extreme-action process unchanged. This is
the action-count analogue of the classical conditioning/extremes
phenomenon discussed in \cite{janson}.

\begin{theorem}[Poisson limit for conditioned large actions]
\label{thm:PPP}
Under the compact critical-window assumptions of
Section~\ref{sec:cutoff}, condition on $S_{n,\beta_n}=n$. The point process
\begin{equation}\label{eq:Xi}
 \Xi_n=\sum_{j\ge1}Z_j\,\delta_{j/e_n}
\end{equation}
converges vaguely on $(0,\infty)$ to a Poisson point process with intensity
$x^{-3}\dd x$. More explicitly, for every nonnegative continuous function
$f$ with compact support in $(0,\infty)$,
\begin{equation}\label{eq:LaplacePPP}
 \E\left[e^{-\int f\dd\Xi_n}\mid S=n\right]
 \longrightarrow
 \exp\left(\int_0^\infty(e^{-f(x)}-1)x^{-3}\dd x\right),
\end{equation}
uniformly for $\lambda_n$ in a fixed compact interval.
For every $s>0$, the conditioned number of actions larger than $se_n$
converges to $\Pois(1/(2s^2))$.
\end{theorem}
\begin{proof}
The joint Fourier transform with the Laplace weight differs from the
unweighted transform by the factor
\begin{equation}\label{eq:PPPfactor}
 \exp\left(n\beta_n\sum_{j\ge1}a_j
  (e^{-f(j/e_n)}-1)e^{ijs/b_n}\right).
\end{equation}
For $s$ in a bounded interval, $j/e_n$ remains in a fixed compact subset
of $(0,\infty)$ wherever $f$ is nonzero. Since $e_n/b_n\to0$,
$e^{ijs/b_n}\to1$ uniformly there. The Riemann sum, using
$n\beta_n a=e_n^2$, tends to
$\int(e^{-f(x)}-1)x^{-3}\dd x$.

The exponent of \eqref{eq:PPPfactor} has uniformly bounded modulus on
the entire Fourier interval: the sum of its absolute coefficients is
$O(n\beta_n\sum_{j\ge c e_n}a_j)=O(1)$, where $c>0$ is below the
support of $f$. Thus the unweighted Fourier damping is preserved up to
a constant factor. Dominated localization as in
Theorem~\ref{thm:allorders} gives the limiting Gaussian integral times
the limiting factor in \eqref{eq:PPPfactor}. Division by
$\Prob(S=n)$ proves \eqref{eq:LaplacePPP}, with uniformity on compact
$\lambda_n$ intervals.

To pass from compactly supported tests to counts on $(s,\infty)$, control
the far tail. The Poisson shift identity gives
\[
 \E\left[\sum_{j>Re_n}Z_j\mid S=n\right]
 =n\beta_n\sum_{j>Re_n}a_j
        \frac{\Prob(S=n-j)}{\Prob(S=n)}.
\]
The uniformly additive local limit bounds the numerator by $C/b_n$
for every $j$, while the denominator is at least $c_0/b_n$ in the
compact window. Hence the displayed expectation is $O(R^{-2})$,
uniformly in large $n$. This proves tightness of the far tail and the
asserted Poisson count limit.
\end{proof}

If $J_{(k)}$ is the $k$th largest action, counting multiplicity and setting
it to zero when fewer than $k$ actions occur, then
\begin{equation}\label{eq:kth}
 \Prob(J_{(k)}\le se_n\mid S=n)
 \longrightarrow e^{-\mu_s}\sum_{r=0}^{k-1}\frac{\mu_s^r}{r!},
 \qquad \mu_s=\frac1{2s^2}.
\end{equation}
The case $k=1$ reproduces \eqref{eq:leadingcutoff}. This derivation also
shows which step fails in the fractional stable regime: there the extreme
and fluctuation scales are comparable, so the Fourier phase of a large
action does not disappear.

\section{Moving finite-action folds: a convergent all-order drift chart}
\label{sec:folds}

At $\beta=\bc$, every finite truncation has an ordinary analytic fold,
even though the limiting critical point is logarithmic. The location of
that fold determines the exponentially small coefficient loss when the
cutoff lies below the extreme scale. It is essential to distinguish its
exact location from a truncated asymptotic approximation.

\subsection{Critical point, critical value, and curvature}
For $M>J$, let $\tau_M>1$ be the unique solution of
\begin{equation}\label{eq:taudef}
 \bc\sum_{j\le M}j a_j\tau_M^j=1.
\end{equation}
Existence and uniqueness follow because the left side is strictly
increasing, is below $1$ at $\tau=1$, and tends to infinity. Define
\begin{equation}\label{eq:folddefs}
 \sigma_M=\log\tau_M,\quad y_M=M\sigma_M,\quad
 r_M=\tau_M e^{-\bc A_M(\tau_M)},\quad
 \mathcal V_M=\bc\sum_{j\le M}j^2a_j\tau_M^j,
\end{equation}
and
\begin{equation}\label{eq:LM}
 L_M=H_M+\frac{p^{[2]}}a.
\end{equation}
The second derivative of $\log q$ with respect to $\log t$ at the
critical point is $-\mathcal V_M<0$, so this is a nondegenerate analytic
fold. The use of $\mathcal V_M$ here avoids confusion with the untilted
quantity $v_{n,M}$ in \eqref{eq:nukappav}.

Introduce entire functions
\begin{align}
 F(y)&=\sum_{k\ge2}\frac{y^k}{k!(k-1)}
      =\int_0^1\frac{e^{yt}-1-yt}{t^2}\dd t,\label{eq:F}\\
 G(y)&=\sum_{k\ge3}\frac{y^k}{k!(k-2)},\qquad G'(y)=F(y),\label{eq:G}\\
 J(y)&=\sum_{k\ge1}\frac{y^k}{k!k}.
 \label{eq:J}
\end{align}
The removable singularities at $t=0$ in the integral representation are
understood by Taylor expansion.

\begin{theorem}[Convergent finite-fold drift functions]
\label{thm:foldchart}
There is a unique function $Y(z)$ analytic near $z=0$ such that
\begin{equation}\label{eq:Y}
 Y=z(1-F(Y)),\qquad Y(0)=0.
\end{equation}
Define the analytic function, with its removable value at $z=0$,
\begin{equation}\label{eq:H}
 \mathcal H(z)=\frac12+Y(z)-\frac{Y(z)^2}{2z}-G(Y(z)).
\end{equation}
Then
\begin{align}
 y_M&=Y(1/L_M)+O((M L_M)^{-1}),\label{eq:yM}\\
 \frac{M^2}{c}\log\frac{r_M}{\rhoc}
     &=\mathcal H(1/L_M)+O(M^{-1}),\label{eq:rM}\\
 \frac{\mathcal V_M}{c}
     &=L_M+J(Y(1/L_M))+O(M^{-1}).\label{eq:VM}
\end{align}
Thus the complete inverse-logarithmic expansions of these quantities
are obtained by substituting $z=1/L_M$ into convergent power series.
\end{theorem}

\begin{proof}
The implicit equation $Y-z(1-F(Y))=0$ has derivative $1$ in $Y$ at
$(Y,z)=(0,0)$, so the analytic implicit function theorem proves the
first assertion. Since $Y(z)=z+O(z^3)$, the quotient in
\eqref{eq:H} is removable and $\mathcal H$ is analytic.

Subtract \eqref{eq:taudef} at $\tau=1$ from its value at $\tau_M$.
The tail is unaltered for $M>J$, so
\begin{equation}\label{eq:foldcriticalexact}
 \sum_{j\le M}j a_j(e^{j y_M/M}-1)
     =a\sum_{j>M}j^{-2}.
\end{equation}
The linear term on the left is $a y_M L_M/M$. For bounded $y\ge0$,
the remaining tail terms, after multiplication by $M/a$, form a
Riemann sum for $F(y)$, with error $O(y^2/M)$; the finite-prefix
nonlinear terms have the same error order. More explicitly, the
function $(e^{yt}-1-yt)/t^2$ extends smoothly to $[0,1]$, and its
supremum and first derivative there are $O(y^2)$ for bounded $y$.
Meanwhile
$M\sum_{j>M}j^{-2}=1+O(M^{-1})$. It follows that
\begin{equation}\label{eq:yapprox}
 L_M y_M+F(y_M)=1+O(M^{-1}).
\end{equation}
Positivity in \eqref{eq:foldcriticalexact} gives $y_M=O(1/L_M)$
(the total linear moment is positive for large $M$).
The exact equation $L_My+F(y)=1$ has the solution $Y(1/L_M)$.
Its derivative on the intervening small positive interval is
$L_M+F'(y)\ge L_M/2$ for large $M$. The mean value theorem proves
\eqref{eq:yM}.

For the critical value, cancel the first-order terms before estimating:
\begin{equation}\label{eq:radius-exact}
 \log\frac{r_M}{\rhoc}
 =\bc\sum_{j>M}a_j
  +\bc\sigma_M\sum_{j>M}j a_j
  -\bc\sum_{j\le M}a_j(e^{j\sigma_M}-1-j\sigma_M).
\end{equation}
Its quadratic term is $-c y_M^2L_M/(2M^2)$. The terms of order at
least three give a Riemann sum for $-cG(y_M)/M^2$, with error
$O(M^{-3})$ for the bounded $y_M$ under consideration. Using
$M^2\sum_{j>M}j^{-3}=1/2+O(M^{-1})$ and
$M\sum_{j>M}j^{-2}=1+O(M^{-1})$, we obtain
\[
 \frac{M^2}{c}\log\frac{r_M}{\rhoc}
 =\frac12+y_M-\frac{L_My_M^2}{2}-G(y_M)+O(M^{-1}).
\]
Replace $y_M$ by $Y(1/L_M)$ using \eqref{eq:yM}. The replacement
error is $O(M^{-1})$; in fact the first derivative of the main
expression vanishes at that exact solution. This proves \eqref{eq:rM}.

Finally,
\[
 \frac{\mathcal V_M}{c}=L_M+
 \sum_{j\le M}\frac{e^{j y_M/M}-1}{j}+O(y_M/M).
\]
The finite-prefix term has been separated, and the sum is a Riemann sum
for $J(y_M)$ with error $O(M^{-1})$. Substituting
\eqref{eq:yM} proves \eqref{eq:VM}.
\end{proof}

\subsection{Coefficients, cancellation, and explicit drift laws}
The first terms of the convergent functions are
\begin{align}
 Y(z)={}&z-\frac12z^3-\frac1{12}z^4+\frac{35}{72}z^5
                +\frac{33}{160}z^6-\frac{1013}{1800}z^7+O(z^8),
 \label{eq:Ycoeff}\\
 \mathcal H(z)={}&\frac12+\frac z2-\frac{z^3}{6}
      -\frac{z^4}{48}+\frac{11z^5}{90}
      +\frac{119z^6}{2880}+O(z^7),\label{eq:Hcoeff}\\
 J(Y(z))={}&z+\frac{z^2}{4}-\frac{4z^3}{9}
      -\frac{31z^4}{96}+\frac{653z^5}{1800}+O(z^6).
 \label{eq:Jcoeff}
\end{align}
They follow by substitution in \eqref{eq:Y} and ordinary power-series
arithmetic. In particular,
\begin{align}
 M L_M\log\tau_M&=1-\frac1{2L_M^2}
      -\frac1{12L_M^3}+O(L_M^{-4}+M^{-1}),\label{eq:tauexp}\\
 \log\frac{r_M}{\rhoc}
 &=\frac{c}{2M^2}\left(1+\frac1{L_M}-\frac1{3L_M^3}
              +O(L_M^{-4}+M^{-1})\right),\label{eq:rexp}\\
 \mathcal V_M&=c\left(L_M+\frac1{L_M}
             +\frac1{4L_M^2}+O(L_M^{-3}+M^{-1})\right).
 \label{eq:Vexpfold}
\end{align}
There is no $L_M^{-2}$ term in the bracket in \eqref{eq:rexp}.
That cancellation is not visible from the leading estimate
$\sigma_M\asymp(M\log M)^{-1}$ alone. The single analytic function
$\mathcal H$ determines every later cancellation and correction.

\begin{remark}[Why convergence here does not imply convergence everywhere]
The series for $Y$, $\mathcal H$, and $J\circ Y$ converge near $0$.
After substituting $1/L_M$, they describe the logarithmic part of the
finite-$M$ expansion, but an additional error of order $M^{-1}$ remains
in the normalized formulas. This is an analytic two-scale expansion,
not an assertion that a power series in $1/\log M$ represents the
finite fold exactly for each $M$.
\end{remark}

\subsection{A uniform saddle law for finite truncations}
\begin{theorem}[Triangular finite-fold coefficient law]
\label{thm:foldLLT}
At $\beta=\bc$, suppose $M=M_n\to\infty$ and
\begin{equation}\label{eq:foldcondition}
 \frac{M_n}{\sqrt{n\log M_n}}\longrightarrow0.
\end{equation}
With the \emph{exact} quantities in \eqref{eq:folddefs},
\begin{equation}\label{eq:foldcoeff}
 \boxed{\quad
 u_n^{(M)}(\bc)\sim
 \frac{r_M^{-n}}{\sqrt{2\pi}\,n^{3/2}\sqrt{\mathcal V_M}}.
 \quad}
\end{equation}
\end{theorem}
\begin{proof}
Exponentially tilt the finite Poisson family to independent variables
with means $n\bc a_j\tau_M^j$, $j\le M$. Its sum $\widetilde S_M$
has mean $n$ by \eqref{eq:taudef} and variance $n\mathcal V_M$.
Lagrange inversion and exponential tilting give exactly
\begin{equation}\label{eq:tiltedidentity}
 u_n^{(M)}(\bc)=\frac{r_M^{-n}}n\Prob(\widetilde S_M=n).
\end{equation}
Set $B_n=\sqrt{n\mathcal V_M}\asymp\sqrt{n\log M}$. Since
$\tau_M^j\le e^{y_M}=1+o(1)$, the third absolute cumulant divided by
$B_n^3$ is at most a constant times
\[
 \frac{nM}{(n\log M)^{3/2}}
 =O\left(\frac{M/B_n}{\log M}\right)\longrightarrow0.
\]
For every fixed Fourier frequency on the $B_n$ scale, the centered
characteristic function therefore tends to $e^{-s^2/2}$.

For completeness, the local limit also requires tail control. The tilt
only increases the positive weights, so
Lemma~\ref{lem:damping} applies. On $|\theta|\le M^{-1/2}$, its
truncated lower bound is at least a constant times
$\theta^2\log M$, giving $e^{-Cs^2}$ on the $B_n$ scale.
On $M^{-1/2}\le|\theta|\le\theta_0$, it is at least a constant
times $n\log M/M$. Under \eqref{eq:foldcondition}, this quantity
dominates $\log B_n$. The remaining interval has damping of order $n$.
Thus dominated Fourier inversion gives
$\Prob(\widetilde S_M=n)\sim(2\pi n\mathcal V_M)^{-1/2}$.
Substitute in \eqref{eq:tiltedidentity}.
\end{proof}

The sufficient condition \eqref{eq:foldcondition} includes $M\asymp
\sqrt n$, where the finite-fold description overlaps with the cutoff
profile. Dividing \eqref{eq:foldcoeff} by the full critical leading
coefficient gives
\begin{equation}\label{eq:matching}
 R_{n,M}\sim
 \exp\left(-n\log\frac{r_M}{\rhoc}\right)
       \sqrt{\frac{c\ell_n}{\mathcal V_M}}.
\end{equation}
At $M=s_ne_n$ with $s_n\to s>0$, \eqref{eq:rexp} yields
$n\log(r_M/\rhoc)\to1/(2s^2)$ and the square-root factor tends to
$1$. Hence the analytic moving-fold calculation agrees with the
probabilistic cutoff law. The refined formula in
Theorem~\ref{thm:refined} provides more precision than the leading
saddle equivalence \eqref{eq:matching}.

\begin{proposition}[When a drift approximation may be exponentiated]
\label{prop:exponentiation}
Let $\mathcal H_K$ be the Taylor polynomial of $\mathcal H$ through
degree $K$, and set
\[
 \Delta_{M,K}=\frac{c}{M^2}\mathcal H_K(1/L_M).
\]
Then
\begin{equation}\label{eq:drifterror}
 \log(r_M/\rhoc)-\Delta_{M,K}
 =O\left(M^{-2}L_M^{-K-1}+M^{-3}\right).
\end{equation}
Replacing $(r_M/\rhoc)^{-n}$ by $e^{-n\Delta_{M,K}}$ incurs a
relative error tending to zero whenever
\begin{equation}\label{eq:expcondition}
 \frac{n}{M^2L_M^{K+1}}+\frac{n}{M^3}\longrightarrow0.
\end{equation}
\end{proposition}
\begin{proof}
The analyticity of $\mathcal H$ gives a Taylor remainder
$O(L_M^{-K-1})$, and \eqref{eq:rM} gives the $O(M^{-3})$ error
after rescaling. Multiply by $n$ before exponentiating.
\end{proof}
If only the leading drift $c/(2M^2)$ is used, its first omitted term is
$c/(2M^2L_M)$, which may become large after multiplication by $n$ in a
hard-cutoff regime. A correct leading statement about $r_M$ is not
automatically a correct relative statement about $r_M^{-n}$.

\section{A slowly varying extension: logarithmically weighted cubic tails}
\label{sec:weighted}

The exact polylogarithm is responsible for the clean all-order analytic
remainder in \eqref{eq:normal}. The leading scale separation needs much
less. This section treats a concrete slowly varying class, without
asserting the same all-order corrections under weaker hypotheses.

Assume now that $a_j\ge0$, $a_1>0$, and
\begin{equation}\label{eq:weightedtail}
 a_j\sim a j^{-3}(\log j)^r\qquad(j\to\infty),
 \qquad a>0,\quad r\ge-1.
\end{equation}
Only large $j$ are involved in this assumption. Define $A(t)=\sum a_jt^j$,
$d_0=A(1)$, $d_1=\sum j a_j$, $\bc=d_1^{-1}$, and
$\rho_\beta=e^{-\beta d_0}$ as before. The parametrization, Lagrange
identity, and Poisson formulas of Section~\ref{sec:model} remain exact.
Put $L_n^*=\tfrac12\log n$, and use the following \emph{leading-order}
scales for $\beta$ in a compact positive interval:
\begin{align}
 e_n^*&=\sqrt{n\beta a}\,(L_n^*)^{r/2},\label{eq:weightedextreme}\\
 b_n^*&=
 \begin{cases}
 \displaystyle\sqrt{\frac{n\beta a}{r+1}}\,
                (L_n^*)^{(r+1)/2},&r>-1,\\[3pt]
 \displaystyle\sqrt{n\beta a\log\log n},&r=-1.
 \end{cases}\label{eq:weightednormal}
\end{align}
These are not exact Lambert scales and should not be substituted into
the correction formulas of Sections~\ref{sec:allorders}--\ref{sec:cutoff}.

\begin{lemma}[Tail and truncated-variance estimates]
\label{lem:weightedtails}
Under \eqref{eq:weightedtail}, as $x\to\infty$,
\begin{align}
 \sum_{j>x}a_j&\sim\frac a2 x^{-2}(\log x)^r,
 &\sum_{j>x}j a_j&\sim a x^{-1}(\log x)^r,\label{eq:weightedtails}\\
 V(x):=\sum_{j\le x}j^2a_j&\sim
 \begin{cases}
 \displaystyle\frac a{r+1}(\log x)^{r+1},&r>-1,\\[3pt]
 a\log\log x,&r=-1.
 \end{cases}\label{eq:weightedvariance}
\end{align}
Furthermore, $\sum_{j\le x}j^3a_j=O(x(\log x)^r)$.
\end{lemma}
\begin{proof}
For each fixed $\epsilon>0$, the weights beyond a fixed index lie
between $(1-\epsilon)$ and $(1+\epsilon)$ times their comparison
sequence. Integral comparison, followed by $\epsilon\downarrow0$,
reduces the assertions to the corresponding integrals. For the tails,
integration by parts shows that differentiating the logarithmic factor
produces a relative $O(1/\log x)$ error. For the variance, integrate
$(\log t)^r/t$; it has primitive $(\log t)^{r+1}/(r+1)$ when
$r>-1$, and $\log\log t$ when $r=-1$. Finally,
$\int^x(\log t)^r\dd t=O(x(\log x)^r)$ for fixed $r$.
Finite prefixes are negligible because $V(x)\to\infty$.
\end{proof}

\begin{theorem}[Log-weighted critical chart at leading order]
\label{thm:weightedchart}
At $\beta=\bc$, the small positive critical inverse satisfies
\begin{equation}\label{eq:weightednormalform}
 \eps=x+\bc(A(e^{-x})-d_0)
       \sim\frac{\bc}{2}x^2V(1/x),\qquad x\downarrow0.
\end{equation}
Consequently,
\begin{equation}\label{eq:weightedinverse}
 x\sim
 \begin{cases}
 \displaystyle
 \sqrt{\frac{2(r+1)\eps}{\bc a}}\,
       \left(\frac12\log\frac1\eps\right)^{-(r+1)/2},&r>-1,\\[6pt]
 \displaystyle
 \sqrt{\frac{2\eps}{\bc a\log\log(1/\eps)}},&r=-1.
 \end{cases}
\end{equation}
In both cases the critical inverse has no nonzero convergent finite
Puiseux representation at $\eps=0$.
\end{theorem}
\begin{proof}
Criticality cancels the linear term exactly:
\[
 \eps=\bc\sum_{j\ge1}a_j(e^{-jx}-1+jx).
\]
For $j\le1/x$, Taylor's formula differs from $j^2x^2/2$ by at most
a constant times $j^3x^3$. Lemma~\ref{lem:weightedtails} bounds the
sum of these errors by $O(x^2(\log(1/x))^r)$. For $j>1/x$,
$0\le e^{-jx}-1+jx\le jx$, so the tail has the same upper bound.
Both errors are $o(x^2V(1/x))$ for every $r\ge-1$. This proves
\eqref{eq:weightednormalform}.

Taking logarithms first gives
$\log(1/x)\sim\tfrac12\log(1/\eps)$. Substituting this into the
slowly varying factors in \eqref{eq:weightedvariance} gives
\eqref{eq:weightedinverse}. In each case the logarithmic order of
$x$ is $\eps^{1/2}$ but $x/\sqrt\eps\to0$, excluding a nonzero
convergent finite Puiseux series exactly as in
Proposition~\ref{prop:ramification}.
\end{proof}

\begin{theorem}[Two scales and a universal leading cutoff law]
\label{thm:weightedbudget}
Let $\beta=\beta_n$ lie in a compact positive interval, and suppose
\[
 \lambda_n^*=\frac{n(1-\beta_n d_1)}{b_n^*}
\]
lies in a compact real interval. Then, uniformly in that window,
\begin{equation}\label{eq:weightedLLT}
 u_n(\beta_n)=\frac{\rho_{\beta_n}^{-n}}{n b_n^*}
                      \bigl(\ph(\lambda_n^*)+o(1)\bigr).
\end{equation}
If $M/e_n^*\to s\in(0,\infty)$, then
\begin{equation}\label{eq:weightedcutoff}
 \frac{u_n^{(M)}(\beta_n)}{u_n(\beta_n)}
       \longrightarrow e^{-1/(2s^2)}.
\end{equation}
More generally, the ratio tends to $1$ if and only if $M/e_n^*\to\infty$,
and tends to zero when $M/e_n^*\to0$. The conditioned extreme-action
point process at scale $e_n^*$ has the same Poisson limit with intensity
$x^{-3}\dd x$ as in Theorem~\ref{thm:PPP}.
\end{theorem}
\begin{proof}
The displayed scales satisfy
\begin{equation}\label{eq:weightedscalerelations}
 \frac{n\beta V(b_n^*)}{(b_n^*)^2}\longrightarrow1,
 \qquad \frac{e_n^*}{b_n^*}\longrightarrow0,
 \qquad n\beta\sum_{j>se_n^*}a_j\longrightarrow\frac1{2s^2}.
\end{equation}
Indeed, $\log b_n^*\sim\log e_n^*\sim\tfrac12\log n$, and
Lemma~\ref{lem:weightedtails} applies. For $r=-1$ the squared ratio
$(b_n^*/e_n^*)^2$ is asymptotic to
$(\tfrac12\log n)\log\log n$; for $r>-1$ it is asymptotic to
$\log n/(2(r+1))$.

For fixed real $s$, split the centered characteristic exponent at
$j=b_n^*$. The part with $j\le b_n^*$ has quadratic term
$-s^2 n\beta V(b_n^*)/(2(b_n^*)^2)\to-s^2/2$. The absolute
Taylor remainder is bounded by a constant times
\[
 \frac{n\beta}{(b_n^*)^3}\sum_{j\le b_n^*}j^3a_j
 =O\left(\frac{(\log b_n^*)^r}{V(b_n^*)}\right)=o(1).
\]
For $j>b_n^*$, use $|e^{iu}-1-iu|\le2+|u|$ and the tail
estimates; the bound has the same order. This proves the Gaussian
characteristic limit.

The local limit follows by the same localization argument as before.
At small $\theta$, summing over $j\le1/|\theta|$ gives the lower
bound
$C\theta^2V(1/|\theta|)$ for the damping. On
$|s|\le(b_n^*)^{1/2}$, the ratio
$V(b_n^*/|s|)/V(b_n^*)$ is bounded below by a positive constant,
using \eqref{eq:weightedvariance}; hence there is Gaussian damping
on that interval. On the next interval the exponent is bounded below
by a positive multiple of $b_n^*$, after choosing a sufficiently small
fixed angular endpoint. This follows from the eventual increase of
$\theta^2V(1/\theta)$ for the explicit logarithmic comparison functions.
Outside it, $a_1>0$ supplies a fixed gap. Dominated Fourier inversion
therefore gives an additive local limit, uniformly at all integer
arguments. Equation~\eqref{eq:poisson} proves \eqref{eq:weightedLLT}.

At $M\sim se_n^*$, the retained variance has the same leading size:
$V(M)/V(b_n^*)\to1$. The omitted mean on the $b_n^*$ scale tends
to zero, since
\[
 \frac{n\beta}{b_n^*}\sum_{j>M}j a_j
       =O(e_n^*/b_n^*)=o(1).
\]
The third scaled cumulant of the truncated family tends to zero as
well. Truncated damping, obtained by stopping the lower-bound sum at
$\min(M,1/|\theta|)$, proves its Gaussian local limit. The exact
identity \eqref{eq:cutoff-exact} and \eqref{eq:weightedscalerelations}
then give \eqref{eq:weightedcutoff}. Monotonicity proves the budget
criterion for arbitrary cutoffs, as in Theorem~\ref{thm:budget}.

Finally, the Riemann sums at scale $e_n^*$ satisfy
$n\beta\sum a_j f(j/e_n^*)\to\int f(x)x^{-3}\dd x$ for compactly
supported continuous $f$ away from zero. The relative tail asymptotic
in \eqref{eq:weightedtail} is uniform over such index ranges once their
lower endpoint tends to infinity. Since $e_n^*/b_n^*\to0$, the
Fourier argument of Theorem~\ref{thm:PPP} applies unchanged. Its
uniform local-limit bound also supplies the necessary far-tail control.
\end{proof}

\begin{remark}[What is not inherited by weakening the tail hypothesis]
An asymptotic relation $a_j\sim a j^{-3}(\log j)^r$ allows arbitrary
slow convergence and oscillatory lower-order errors. It does not imply
an analytic remainder, a convergent Lambert--analytic chart, or the
rate in Theorem~\ref{thm:refined}. Theorems
\ref{thm:weightedchart} and \ref{thm:weightedbudget} deliberately make
only leading-order assertions. The case $r=-1$ also shows that simply
replacing a divergent variance by a single $\log n$ is not valid
throughout the slowly varying class.
\end{remark}

\section{Reproducible computation and numerical diagnostics}
\label{sec:computation}

\subsection{Exact checks and numerical method}
The accompanying program \texttt{code/verify.py} is a standalone
reproduction script. Its exact rational/symbolic checks verify the
Lambert chart equation through $T^2$, the finite-fold series $Y$ through
$z^7$, and the Lagrange coefficient identity through degree $6$ in a
rational test model. The additional drift and curvature series are
checked by exact substitution. These finite checks detect algebraic
mistakes; they are not substitutes for the all-order proofs.

For numerical diagnostics we use the pure model $a=1$, $P=0$, so
\[
 d_0=\zeta(3),\quad d_1=\zeta(2),\quad \bc=1/\zeta(2),\quad h=3/2.
\]
The script computes normalized coefficients with a discrete Fourier
coefficient extraction at radius $r=e^{-3/n}$, using a power-of-two
period $L\ge16(n+1)$. Terms in $A$ above degree $n$ can be discarded
\emph{exactly} when extracting the coefficient of degree $n$.
The normalized exponential has nonnegative coefficients with total
mass at most $1$. Therefore its coefficient aliases, in exact arithmetic,
are nonnegative and have total error at most $r^L=e^{-3L/n}$ after
dividing by $r^n$. Floating-point roundoff is separate from that alias
bound and is not rigorously enclosed.

An independent 80-decimal-digit coefficient recurrence at $n=128$
agreed with the FFT value to a relative discrepancy of about
% ed. (2026-09-29): 6.76 corrected to 6.75; the recorded value is 6.75228185529e-15 (data/verification_run.txt, BUILD_REPORT.json).
$6.75\times10^{-15}$. The recurrence is
\begin{equation}\label{eq:coeffrecurrence}
 p_0=e^{-n\beta d_0},\qquad
 p_k=\frac{n\beta}{k}\sum_{j=1}^{\min(k,M)}j a_jp_{k-j},
\end{equation}
with $M=n$ for the full coefficient. No conclusion about global
floating-point certification follows from this one comparison.

\subsection{Sector coefficients: inverse logarithms converge slowly}
Write $p_n=\Prob(S_{n,\bc}=n)$. Table~\ref{tab:coefficients}
compares the normalized mass with the first two corrections in
Corollary~\ref{cor:criticalcoeff}.
\begin{table}[htbp]
\centering\small
\begin{tabular}{rrrrr}
\toprule
$n$ & $\ell_n$ & $b_n p_n/\phi(0)$ & $1+g_1/\ell_n$ & $1+g_1/\ell_n+g_2/\ell_n^2$ \\
\midrule
256 & 4.8090 & 0.98892962 & 1.03793096 & 1.00781273 \\
1024 & 5.5761 & 0.99803178 & 1.03271250 & 1.01031140 \\
4096 & 6.3329 & 1.00306096 & 1.02880338 & 1.01143623 \\
16384 & 7.0820 & 1.00592159 & 1.02575691 & 1.01186924 \\
65536 & 7.8250 & 1.00756641 & 1.02331113 & 1.01193568 \\
\bottomrule
\end{tabular}

\caption{Critical sector diagnostics for the pure cubic-tail model. The
last two columns are asymptotic approximations, not upper bounds.}
\label{tab:coefficients}
\end{table}
Even at $n=65536$, $\ell_n$ is only about $7.825$. The two-correction
approximation is closer in these examples, but convergence is visibly
slow. This behavior is consistent with an inverse-logarithmic expansion
and warns against estimating its accuracy from the size of $n$ alone.

\subsection{Cutoff and critical-window diagnostics}
Table~\ref{tab:cutoffs} uses $M$ nearest to $e_n$; the effective value
$s_n=M/e_n$ is shown to avoid silently replacing an integer cutoff by
its continuous target. The refined prediction exponentiates the
explicit part of \eqref{eq:refinedexact}, retaining the exact finite
sums $\nu_{n,M}$, $\kappa_{n,M}$, and $v_{n,M}$.
\begin{table}[htbp]
\centering\small
\begin{tabular}{rrrrrr}
\toprule
$n$ & $M$ & $s_n$ & Actual ratio & Leading & Refined \\
\midrule
256 & 12 & 0.9619 & 0.6241867 & 0.5825271 & 0.6209904 \\
1024 & 25 & 1.0020 & 0.6424435 & 0.6077365 & 0.6395692 \\
4096 & 50 & 1.0020 & 0.6371199 & 0.6077365 & 0.6347475 \\
16384 & 100 & 1.0020 & 0.6344374 & 0.6077365 & 0.6324177 \\
65536 & 200 & 1.0020 & 0.6329721 & 0.6077365 & 0.6312248 \\
\bottomrule
\end{tabular}

\caption{Retained coefficient fractions near the critical action budget.
The refined prediction incorporates omitted mass, mean, and variance.}
\label{tab:cutoffs}
\end{table}
The leading prediction approaches the correct limit but misses
substantial finite-$n$ corrections. The refined formula reduces those
discrepancies in the displayed range; the remaining differences are
not treated as certified error intervals.

Table~\ref{tab:coupling} changes $\beta$ at $n=4096$ to obtain
$\lambda_n=-1,0,1$. The density ratio is
$b_n\Prob(S=n)/\ph(\lambda_n)$, its first correction is
$1+g(\lambda_n)/\ell_n$, and the cutoff is again nearest to $e_n$.
The cutoff's leading limit is independent of $\lambda_n$, but its
finite-window correction is not.
\begin{table}[htbp]
\centering\small
\begin{tabular}{rrrrr}
\toprule
$\lambda_n$ & Density ratio & First correction & Cutoff ratio & Refined cutoff \\
\midrule
-1 & 1.1651828 & 1.1303772 & 0.8218048 & 0.8143599 \\
0 & 1.0030610 & 1.0288034 & 0.6371199 & 0.6347475 \\
1 & 0.8072274 & 0.8039156 & 0.3716506 & 0.3690131 \\
\bottomrule
\end{tabular}

\caption{Finite-coupling diagnostics at sector index $4096$.}
\label{tab:coupling}
\end{table}

\subsection{Inverse chart and moving folds}
In the same pure model, the relative errors of the core $T$ and its
first two analytic corrections are as follows:
\begin{center}
\small
\begin{tabular}{rrrr}
\toprule
$\eps$ & Core $T$ & $T+v_1T^2$ & $T+v_1T^2+v_2T^3$\\
\midrule
$10^{-4}$ & $1.0020\cdot10^{-4}$ & $5.5798\cdot10^{-8}$ & $2.7144\cdot10^{-11}$\\
$10^{-8}$ & $4.1471\cdot10^{-7}$ & $1.3691\cdot10^{-12}$ & $2.9478\cdot10^{-18}$\\
$10^{-16}$ & $1.6274\cdot10^{-11}$ & $3.3625\cdot10^{-21}$ & $3.0045\cdot10^{-31}$\\
\bottomrule
\end{tabular}
\end{center}
The reference roots were computed at 80-decimal-digit working precision,
not enclosed by interval arithmetic. Here the exact Lambert core has
already absorbed the logarithmic singularity, and the remaining
analytic expansion is correspondingly effective.

For finite folds, Table~\ref{tab:folds} compares the exact numerical
critical point and value with \eqref{eq:tauexp} and \eqref{eq:rexp}.
The radius column tends to $1$, while the three-term prediction includes
the $L_M^{-1}$ and $L_M^{-3}$ terms. Radius drift was evaluated using
the cancellation-resistant identity \eqref{eq:radius-exact}.
\begin{table}[htbp]
\centering\small
\begin{tabular}{rrrr}
\toprule
$M$ & $M L_M\log\tau_M$ & $2M^2\log(r_M/\rho_c)/c$ & Three-term prediction \\
\midrule
32 & 0.95564674 & 1.20338010 & 1.24141038 \\
128 & 0.97935090 & 1.17279266 & 1.18197701 \\
512 & 0.98825384 & 1.14341348 & 1.14565007 \\
2048 & 0.99228716 & 1.12076653 & 1.12131621 \\
8192 & 0.99446566 & 1.10378025 & 1.10391682 \\
\bottomrule
\end{tabular}

\caption{Finite-action critical-point and critical-value drift. In this
pure model $L_M=H_M$.}
\label{tab:folds}
\end{table}

\subsection{Reproduction and reliability boundary}
From the package directory, run
\begin{verbatim}
python -m pip install -r requirements.txt
python code/verify.py --max-n 65536
pdflatex -interaction=nonstopmode -halt-on-error article.tex
pdflatex -interaction=nonstopmode -halt-on-error article.tex
\end{verbatim}
A third LaTeX pass may be used to stabilize references. The generated
TeX tables are included, so the PDF can be rebuilt without rerunning the
numerical code. The script resolves output paths relative to its own
package root, not the caller's working directory. It overwrites only
its own generated data files.

The full machine-readable record is \texttt{data/verification.json};
the supplied console output is \texttt{data/verification\_run.txt}.
Exact checks, numerical diagnostics, analytic theorems, and prospective
formalization are distinct layers. The program does not prove Fourier
dominating estimates by sampling, and it does not verify analytic
continuation or quantify a theorem's unstated constants.

\section{Consequences for transseries and proof architecture}
\label{sec:interpretation}

\subsection{Three expansions, three different meanings}
There are three variables and three limiting processes. Near $q=0$,
$U_\beta(q)$ is an ordinary convergent series; after $q=e^{-X}$ it is
an integer-action exponential transseries. Near the critical value
$q=\rhoc$, the local variable is $\eps=\log(\rhoc/q)$, and the
inverse uses the convergent Lambert chart $x=TV(T,\eta)$. Finally,
large sector index $n$ has an asymptotic expansion in $1/\ell_n$,
with exponentially weighted prefactor $\rho_\beta^{-n}$.

Convergence of the small-$q$ germ says nothing by itself about analytic
ramification at $\rhoc$. Convergence of $V$ in two independent
variables does not imply convergence of the large-$n$ inverse-logarithmic
series. The two uses of Lambert $W$ are related by the same logarithmic
balance, but their arguments have different roles: one inverts the
critical map, and the other normalizes a Fourier saddle.

\subsection{What finite-prefix universality actually preserves}
Changing finitely many $a_j$ can change $d_0$, $d_1$, and $h$.
Consequently it changes the critical coupling, the exponential sector
factor, the exact normalization, and the analytic coefficients of the
critical inverse. It does not change the universal functions $G_j$
after that exact normalization, the leading cutoff profile, or the
theorems' logarithmic orders. In the finite-fold chart it is absorbed
in $L_M=H_M+p^{[2]}/a$, up to the displayed algebraically small
errors. These are precise invariance statements, not a claim that
different finite prefixes produce identical functions.

The leading sector expansion also cannot detect every piece of the
analytic remainder. At every fixed inverse-logarithmic order, the
$n/b_n^3$ contribution is smaller than the error scale; it becomes
relevant in a finer multi-scale expansion. This supplies a natural
next level of the transseries rather than a reason to discard the
analytic remainder.

\subsection{A feasible formalization decomposition}
The repository's existing asymptotic-scale and residual-transport work
suggests a modular proof path, but no file in this package is claimed
to be machine-checked. A practical decomposition would be:
\begin{enumerate}[label=\arabic*.]
\item Formalize the exact coefficient identities and the finite-series
recurrences first. These are algebraic and can be separated from
analytic convergence and probability.
\item Formalize the Lambert core identity and the local implicit-chart
interface with explicit hypotheses for holomorphicity, branch choice,
and nonzero derivative. The geometric Cauchy bound is then a reusable
separate theorem.
\item Formalize the finite-fold rescaling identities and Riemann-sum
error estimates, keeping exact critical quantities separate from
asymptotic approximations. Proposition~\ref{prop:exponentiation}
is a particularly useful safety interface.
\item Package Fourier inversion, the damping lemmas, and an
integrable-majorant expansion theorem. The uniform-in-$k$ additive
bound should be represented explicitly, rather than inferred from
pointwise convergence.
\end{enumerate}
The conditioned Poisson process and slowly varying extension could then
be built on those analytic foundations. A declaration-level crosswalk
to existing Lean files would require additional source inspection and
actual compilation; neither is asserted here.

\section{Further research questions and concrete targets}
\label{sec:questions}

\paragraph{1. A genuine two-parameter upper-endpoint crossover.}
Let $a_j=a j^{-1-\alpha_n}$ with $\alpha_n\uparrow2$. Determine the
joint limit when $(2-\alpha_n)\log n$ stays bounded. The truncated
variance involves the cancellation in
$(b^{2-\alpha}-1)/(2-\alpha)$; this suggests a transition parameter,
but neither the critical chart nor the cutoff kernel follows by
substitution. The target is a uniform chart and a uniform coefficient
law connecting the fractional stable regime to this fixed endpoint.
% ed. (2026-09-29): answered by two later sibling packages (batch 49); recorded on filing.
\begin{ednote}
Answered, pending review, by two later packages of this series written
independently of each other. The stable--Gaussian package \cite{ed:sge}
works in exactly this window, $|2-\alpha_n|\log n\le K$, from both sides,
and gives a two-term uniform inverse, a convergent chart about a
desingularized core, the uniform coefficient law with its first
correction, and a uniform first cutoff correction without finite-prefix
moments; it did not see this article. The confluent package \cite{ed:cct}
allows $2-\alpha_n\to0$ at any rate and of either sign, with one
convergent chart in three variables and a leading-order coefficient law;
it read this article but does not cite this question by number. Where both
apply they agree term by term, and at $\alpha=2$ the stable--Gaussian
first cutoff correction is \eqref{eq:criticalcutoffcorrection}, up to the
change of logarithmic normalization.
\end{ednote}

\paragraph{2. Explicit finite-$n$ constants for action selection.}
Turn Theorem~\ref{thm:refined} into an inequality with numerical
constants, threshold values, and interval-enclosed $G_1$. Given $n$,
$\beta$, and a relative loss tolerance, return the smallest certified
cutoff or a certified near-minimal interval. Formula
\eqref{eq:selection} is an asymptotic design rule, not that algorithm.
% ed. (2026-09-29): partial progress in the sibling package filed beside this one (batch 48).
\begin{ednote}
Partial progress: for the same model the marginal package \cite{ed:mct}
proves an explicit finite-$n$ lower bound on the retained fraction (its
\texttt{thm:finite-certificate}) and a two-term minimal budget (its
\texttt{thm:min-budget}). Its recorded evaluations use no directed rounding,
and the certified selection algorithm remains open there too.
\end{ednote}

\paragraph{3. Large orders of the Gaussian--logarithmic kernels.}
Determine the growth of $G_j(\lambda)$ as $j\to\infty$, uniformly on
specified real or complex $\lambda$ sets. Does the inverse-logarithmic
series admit a useful Borel or accelerated summation? Identify an
optimal truncation index and compare its remainder with the first
analytic $n^{-1/2}$-scale corrections. No convergence assertion for
this series is proved above.

\paragraph{4. The finer transseries beyond every inverse-logarithmic order.}
Include the analytic terms of $R_\beta$ in \eqref{eq:char-normal} and
construct a joint expansion in $1/\ell_n$ and $n/b_n^3$. The leading
new Fourier factor is cubic and depends on $r_0$. A desirable result
is an all-order multi-index formula with an error estimate that does
not double-count corrections absorbed in the exact Lambert scale.

\paragraph{5. An expanding critical coupling window.}
Find the largest growing range of $\lambda_n$ on which a finite
correction to \eqref{eq:sectorexp} is relatively accurate. The present
theorem is uniform additively for all integers but only yields a useful
relative law on compact windows. Match large positive and negative
$\lambda_n$ to inherited-tail and interior-saddle regimes, respectively.

\paragraph{6. Beyond logarithmic powers in the slowly varying tail.}
For $a_j\sim a j^{-3}\mathcal L(j)$ with a general slowly varying
$\mathcal L$, characterize conditions under which the extreme and
fluctuation scales separate and conditioning remains asymptotically
invisible to extremes. Specify an effective inverse notation when the
normalization requires a de Bruijn conjugate rather than Lambert $W$.
The explicit log-power subclass does not resolve this general problem.

\paragraph{7. Finite-$M$ corrections to the drift chart.}
Augment Theorem~\ref{thm:foldchart} by powers of $1/M$ with coefficient
functions analytic in $1/L_M$. The tail sums and Riemann sums have
Euler--Maclaurin expansions, but one must keep their finite-prefix
terms and cancellations synchronized. The target is an exponentiable
formula in regimes where $n/M^3$ no longer tends to zero.

\paragraph{8. The lower endpoint and loss of the mean.}
At a tail near $j^{-2}$, $d_1$ may diverge and the critical coupling
$1/d_1$ collapses. Determine an appropriate renormalized coupling,
critical coordinate, and action budget. This is not the same
infinite-variance Gaussian problem and needs a new normalization.

\paragraph{9. Nonlinear slopes and genuinely noncommensurate actions.}
Replace $jU$ by $(j+r_j)U$, or replace integer actions by a sequence
with no common lattice. Determine which endpoint and cutoff laws
persist. The exact substitution $t=qe^U$ or lattice Fourier inversion
then fails, so a stability theorem must replace the exact formulas.

\paragraph{10. Signed weights and competing complex critical values.}
Allow controlled signed or complex perturbations of the action tail.
Find analytic hypotheses replacing positivity and Poisson conditioning,
and identify when multiple complex saddles contribute. Neither the
monotone cutoff criterion nor its probabilistic proof extends
unchanged to that setting.

\paragraph{11. Minimal transseries support and monodromy of the Lambert chart.}
Determine a minimal convenient well-based scale for expanding
$T V(T,\eta)$ after replacing $W_{-1}$ by its nested-logarithmic
expansion. Describe analytic continuation between its logarithmic
branches, with domain bounds. The no-Puiseux proposition rules out only
finite algebraic ramification; it does not classify all Hahn or
logarithmic representations.

\paragraph{12. Proof-producing computation.}
Build a checker for the exact chart recurrences, the finite-fold
stationarity cancellation, and explicit local holomorphic majorants.
Then add interval Fourier quadrature with a proved tail enclosure.
Connect these certificates to formal statements distinguishing exact
coefficient identities, analytic existence, and numerical enclosures.

\section{Conclusion}
The upper endpoint of the positive power-tail action model is governed
by a logarithmic critical singularity. Its local inverse becomes an
analytic convergent function after an exact Lambert change of scale.
Its large-sector coefficients have an independently derived,
arbitrary-order inverse-logarithmic Gaussian expansion. The action
cutoff is controlled by a smaller extreme scale, and its first
finite-window corrections retain information that the leading
Poisson profile suppresses.

The finite-action folds provide a second, analytic explanation of the
same cutoff phenomenon. Their drift and curvature are controlled to
all inverse-logarithmic orders by convergent implicit functions, with
separate algebraic errors that must be multiplied by $n$ before
exponentiation. The log-weighted extension demonstrates that the scale
separation persists beyond the exact polylogarithm, while also making
clear which convergence and rate claims depend on the stronger tail
hypothesis.

These results resolve the fixed-endpoint problem and several defined
extensions within the stated model. They do not settle the full joint
exponent crossover, general resurgence, arbitrary slowly varying
realization, or formal verification. Those remaining targets are now
separated from the statements proved here.

\appendix
\section{Source audit, dependencies, and research status}
\label{app:audit}

\subsection{Repository snapshot and inspected scope}
% ed. (2026-09-29): "pinned tree" corrected to "pinned commit"; the identifier below is a commit.
The repository was accessed at the pinned commit
\begin{center}
\small\texttt{db68f0853c3c69cd930caedab6f6ad9addb11eaf}.
\end{center}
The focused inspection included the directory tree under
\texttt{Analysis/Transseries}, the group README under
\texttt{docs/series-and-transseries}, and the moving-fold package's
README. The group README distinguishes consolidated material from
unmerged research arrivals and records formalization limitations.
This article does not convert those human manuscript claims into
machine-checked results.

The preceding countable-action article \cite{previous} was inspected
as its editable TeX in the user's Library, named
\texttt{article(20260929-193550).tex}. Its model, principal theorem
statements, and research-question section were read. In particular,
its Section~11, Question~4 asks for endpoint logarithmic charts and
cutoff laws, and Questions~1, 2, and 7 motivate the extensions above.
Its presence in the Library does not establish that it was merged into
% ed. (2026-09-29): "tree" corrected to "snapshot" (the pin is a commit).
the pinned GitHub snapshot. This provenance is kept separate from the
repository snapshot.

This was a targeted research audit, not a line-by-line review of the
entire canonical transseries volume, all recent articles, or the Lean
corpus. The proofs in the present article do not assume the correctness
of a new theorem in an unreviewed repository package.

\subsection{External sources and priority boundary}
The primary sources consulted were the NIST DLMF entries for
polylogarithms and Lambert $W$, the author-maintained site for
\emph{Analytic Combinatorics}, Janson's article (including its
Example~18.29 in the arXiv HTML version), and the
Chakrabarty--Samorodnitsky truncation paper's bibliographic record.
The classical leading Gaussian/extreme mechanism is explicitly
acknowledged in Section~\ref{sec:intro}.

The specialized formulas and proofs assembled here warrant independent
mathematical review. A literature search finding no exact counterpart
would still not establish priority by itself. Accordingly, this
package uses ``contribution developed here'' rather than an unsupported
claim that a longstanding worldwide conjecture has been resolved.

\subsection{Proof dependencies}
The model identities in Section~\ref{sec:model} feed two independent
routes. The local analytic route uses the integer-order polylogarithm
expansion and the implicit function theorem to obtain the Lambert
chart. The coefficient route uses the exact Poisson identity,
Fourier inversion, and the damping lemma to obtain all-order sector
estimates. The refined cutoff theorem combines that full local limit
with a separate quantitative truncated local limit. The moving-fold
route is based on exact critical-point identities and controlled
Riemann sums, and only then uses a tilted local limit for coefficients.
The log-weighted extension uses its own leading tail estimates and
does not invoke an analytic remainder it has not established.

\section{Notation ledger}
\small
\begin{longtable}{p{0.20\textwidth}p{0.71\textwidth}}
\toprule
Symbol & Meaning\\
\midrule
\endfirsthead
\toprule
Symbol & Meaning\\
\midrule
\endhead
$A$, $A_M$ & Full action generating function and its cutoff at action $M$.\\
$a_j$, $a$, $P$ & Nonnegative action weights, cubic-tail amplitude, finite-prefix polynomial.\\
$p^{[k]}$ & Finite-prefix moment $\sum_{j\le J}j^kp_j$, not a full moment.\\
$d_0$, $d_1$ & Finite endpoint mass and mean: $A(1)$ and $\sum j a_j$.\\
$\bc$, $\rho_\beta$, $c_\beta$ & Critical coupling $d_1^{-1}$, factor $e^{-\beta d_0}$, and $\beta a$.\\
$h$ & Exact quadratic normal-form constant $3/2+p^{[2]}/a$.\\
$U_\beta$, $u_n$ & Feedback solution and its coefficient of $q^n$.\\
$t$, $x$, $\eps$ & $qe^U$, $-\log t$, and $\log(\rho_\beta/q)$.\\
$R_\beta$, $r_k$ & Analytic normal-form remainder and its critical Taylor coefficients.\\
$T$, $\eta$, $V$ & Exact Lambert core, reciprocal logarithmic coordinate, and convergent analytic chart.\\
$Z_j$, $S$, $S_M$ & Independent Poisson action counts and their full/truncated weighted sums.\\
$e_n$, $b_n$, $\ell_n$ & Extreme scale $\sqrt{n\beta a}$, fluctuation scale, and exact Lambert logarithm.\\
$\lambda_n$, $s_n$ & Critical-window coordinate $n(1-\beta d_1)/b_n$ and cutoff ratio $M/e_n$.\\
$G_j$, $g_j$, $g$ & Fourier correction kernels, their normalized values at $0$, and $G_1/\ph$.\\
$R_{n,M}$ & Retained coefficient fraction $u_n^{(M)}/u_n$.\\
$\nu_{n,M}$, $\kappa_{n,M}$, $v_{n,M}$ & Omitted intensity, omitted scaled mean, and retained scaled variance.\\
$\tau_M$, $r_M$, $\mathcal V_M$ & Exact finite-fold critical point, critical value, and tilted curvature.\\
$\sigma_M$, $y_M$, $L_M$ & $\log\tau_M$, $M\log\tau_M$, and $H_M+p^{[2]}/a$.\\
$F$, $G$, $J$, $Y$, $\mathcal H$ & Entire or local analytic functions describing the finite-fold drift. Here $G$ without a subscript is not the Fourier kernel $G_j$.\\
$e_n^*$, $b_n^*$, $V(x)$ & Leading scales and truncated second moment for log-weighted tails.\\
\bottomrule
\end{longtable}
\normalsize

\begin{thebibliography}{9}

\bibitem{dlmfpoly}
NIST Digital Library of Mathematical Functions.
\emph{Section 25.12: Polylogarithms}, especially Eq.~25.12.12 and its
restriction on the order parameter.
\url{https://dlmf.nist.gov/25.12}. Consulted 29 September 2026.

\bibitem{dlmfw}
NIST Digital Library of Mathematical Functions.
\emph{Section 4.13: Lambert $W$-Function}.
\url{https://dlmf.nist.gov/4.13}. Consulted 29 September 2026.

\bibitem{flajolet}
P. Flajolet and R. Sedgewick.
\emph{Analytic Combinatorics}. Cambridge University Press, 2009.
Author-maintained book site: \url{https://ac.cs.princeton.edu/home/}.

\bibitem{janson}
S. Janson.
Simply generated trees, conditioned Galton--Watson trees, random allocations
and condensation. \emph{Probability Surveys} \textbf{9} (2012), 103--252.
doi:10.1214/11-PS188. arXiv:1112.0510.
The cubic-tail comparison discussed here is Example~18.29 in the inspected
arXiv HTML version: \url{https://arxiv.org/html/1112.0510v1}.

\bibitem{chakrabarty}
A. Chakrabarty and G. Samorodnitsky.
Understanding heavy tails in a bounded world or, is a truncated heavy tail
heavy or not? \emph{Stochastic Models} \textbf{28}(1) (2012), 109--143.
doi:10.1080/15326349.2012.646551. arXiv:1001.3218.

\bibitem{repo}
VladimirReshetnikov/ProveIt.
\emph{Series and transseries}, group README and package inventory.
% ed. (2026-09-29): "Pinned tree" corrected to "Pinned commit".
Pinned commit and audit boundary: Appendix~\ref{app:audit}.
\url{https://github.com/VladimirReshetnikov/ProveIt/tree/db68f0853c3c69cd930caedab6f6ad9addb11eaf/Analysis/Transseries/docs/series-and-transseries}.

\bibitem{fold}
\repo{} repository, research package.
\emph{Critical Transseries at a Moving Fold: A uniform two-sheet atlas,
large-sector crossover, and resonant action splitting}. September 2026.
% ed. (2026-09-29): "pinned tree" corrected to "pinned commit".
Package README inspected at the pinned commit, subdirectory
\texttt{Critical\_Transseries\_Moving\_Fold} of \cite{repo}.

\bibitem{previous}
\emph{Beyond Finite-Action Folds: Critical Hahn Transseries, Stable Sector
Asymptotics, and Sharp Action Budgets}. Earlier research draft prepared
for Vladimir Reshetnikov, 29 September 2026. Editable source inspected
in the user's Library as \texttt{article(20260929-193550).tex};
see especially Section~11, Questions~1, 2, 4, and 7.
Repository inclusion at the pinned snapshot is not assumed.
% ed. (2026-09-29): editorial note on the filed location of the preceding article.
Editorial note (ProveIt, 2026-09-29): the article is now filed as
\path{Analysis/Transseries/docs/series-and-transseries/Critical_Hahn_Transseries_Beyond_Finite_Action_Folds/article.tex}
(batch~47).

% ed. (2026-09-29): the three entries below are editorial additions for the editorial notes.
\bibitem{ed:mct}
[Editorial addition, ProveIt, 2026-09-29.]
\emph{Marginal Critical Transseries: Convergent Lambert-W Charts, Universal
Cutoff Corrections, and Certified Action Budgets}, research package filed
in batch~48 in the group of \cite{repo}:
\path{Marginal_Critical_Transseries_Action_Budgets/Marginal_Critical_Transseries.tex}.

\bibitem{ed:cct}
[Editorial addition, ProveIt, 2026-09-29.]
\emph{Confluent Critical Transseries Across the Exponent-Two Boundary},
research package filed in batch~49 in the group of \cite{repo}:
\path{Confluent_Critical_Transseries_Exponent_Two_Boundary/Confluent_Critical_Transseries.tex}.

\bibitem{ed:sge}
[Editorial addition, ProveIt, 2026-09-29.]
\emph{Through the Stable--Gaussian Endpoint: Uniform Critical Transseries,
Prefix-Independent Cutoff Corrections, and Conditional Extremes}, research
package filed in batch~49 in the group of \cite{repo}:
\path{Stable_Gaussian_Endpoint_Uniform_Critical_Transseries/article.tex}.

\end{thebibliography}
\end{document}
